\documentclass[12pt]{article}

\usepackage{graphicx}
\usepackage{indentfirst}
\usepackage{tocloft}
\usepackage[utf8]{inputenc}
\usepackage{vietnam}
\usepackage{hyphenat}
\usepackage{hyperref}
\usepackage{float}
\usepackage[vietnamese]{babel}
\usepackage{amssymb}
\usepackage{fancyhdr}
\usepackage{xcolor}
\usepackage{datetime}
\usepackage{array}
\usepackage{tikz}
\usepackage{scrextend}
\usetikzlibrary{calc}

\definecolor{myblue}{RGB}{0,0,255}

\hypersetup{
	colorlinks=true,
	linkcolor=myblue,
	citecolor=myblue,
	urlcolor=myblue
}

\addto\captionsvietnamese{\renewcommand{\abstractname}{Mở đầu}}
\usepackage[letterpaper,top=2.5cm,bottom=3cm,left=3.5cm,right=2cm,marginparwidth=1.75cm]{geometry}

\setlength{\parindent}{1.5cm}
\setlength{\parskip}{0.25em}

% Định dạng mục lục phân cấp đẹp
\usepackage{tocloft}

% Thụt lề theo cấp
\setlength{\cftsubsecindent}{1.5em}
\setlength{\cftsubsubsecindent}{3.5em}

% Chấm dẫn cho tất cả các cấp
\renewcommand{\cftsubsecdotsep}{\cftdotsep}
\renewcommand{\cftsubsubsecdotsep}{\cftdotsep}

% Cỡ chữ và kiểu chữ theo cấp
\renewcommand{\cftsecfont}{\bfseries}          % section: in đậm
\renewcommand{\cftsubsecfont}{\normalfont}     % subsection: thường
\renewcommand{\cftsubsubsecfont}{\normalfont\small} % subsubsection: nhỏ hơn




\begin{document}
	%===========================
	% Trang bìa
	%===========================
	\begin{titlepage}
		\begin{tikzpicture}[remember picture,overlay,inner sep=0,outer sep=0]
			\draw[blue!70!black,line width=4pt] ([xshift=-1.5cm,yshift=-2cm]current page.north east) coordinate (A)--([xshift=1.5cm,yshift=-2cm]current page.north west) coordinate(B)--([xshift=1.5cm,yshift=2cm]current page.south west) coordinate (C)--([xshift=-1.5cm,yshift=2cm]current page.south east) coordinate(D)--cycle;
			
			\draw ([yshift=0.5cm,xshift=-0.5cm]A)-- ([yshift=0.5cm,xshift=0.5cm]B)--
			([yshift=-0.5cm,xshift=0.5cm]B) --([yshift=-0.5cm,xshift=-0.5cm]B)--([yshift=0.5cm,xshift=-0.5cm]C)--([yshift=0.5cm,xshift=0.5cm]C)--([yshift=-0.5cm,xshift=0.5cm]C)-- ([yshift=-0.5cm,xshift=-0.5cm]D)--([yshift=0.5cm,xshift=-0.5cm]D)--([yshift=0.5cm,xshift=0.5cm]D)--([yshift=-0.5cm,xshift=0.5cm]A)--([yshift=-0.5cm,xshift=-0.5cm]A)--([yshift=0.5cm,xshift=-0.5cm]A);
			
			\draw ([yshift=-0.3cm,xshift=0.3cm]A)-- ([yshift=-0.3cm,xshift=-0.3cm]B)--
			([yshift=0.3cm,xshift=-0.3cm]B) --([yshift=0.3cm,xshift=0.3cm]B)--([yshift=-0.3cm,xshift=0.3cm]C)--([yshift=-0.3cm,xshift=-0.3cm]C)--([yshift=0.3cm,xshift=-0.3cm]C)-- ([yshift=0.3cm,xshift=0.3cm]D)--([yshift=-0.3cm,xshift=0.3cm]D)--([yshift=-0.3cm,xshift=-0.3cm]D)--([yshift=0.3cm,xshift=-0.3cm]A)--([yshift=0.3cm,xshift=0.3cm]A)--([yshift=-0.3cm,xshift=0.3cm]A);
		\end{tikzpicture}
		
		\begin{center}
			\vspace{7pt}
			BỘ GIÁO DỤC VÀ ĐẠO TẠO
			
			\vspace{7pt}
			\textbf{TRƯỜNG ĐẠI HỌC TÂY BẮC}
		\end{center}
		
		\begin{center}
			\includegraphics[scale=0.5]{images/tb.jpg}
			
			\vspace{10pt}
			\fontsize{18pt}{17pt}\selectfont 
			\textbf{ ĐỒ ÁN TỐT NGHIỆP}
			
			\vspace{10pt}
			% \textbf{\footnotesize MÔN HỌC}
			\vspace{50pt}
		\end{center}
		
		
		\begin{center}
			\fontsize{18pt}{17pt}\selectfont
			\textbf{\textrm{XÂY DỰNG ỨNG DỤNG QUẢN LÝ KÝ TÚC XÁ TRƯỜNG ĐẠI HỌC TÂY BẮC}}
		\end{center}
		
		\vspace{250pt}
		\begin{center}
			Sơn La - năm 2026
		\end{center}
	\end{titlepage}
	
	%===========================
	% Trang bìa
	%===========================
	\begin{titlepage}
		\begin{tikzpicture}[remember picture,overlay,inner sep=0,outer sep=0]
			\draw[blue!70!black,line width=4pt] ([xshift=-1.5cm,yshift=-2cm]current page.north east) coordinate (A)--([xshift=1.5cm,yshift=-2cm]current page.north west) coordinate(B)--([xshift=1.5cm,yshift=2cm]current page.south west) coordinate (C)--([xshift=-1.5cm,yshift=2cm]current page.south east) coordinate(D)--cycle;
			
			\draw ([yshift=0.5cm,xshift=-0.5cm]A)-- ([yshift=0.5cm,xshift=0.5cm]B)--
			([yshift=-0.5cm,xshift=0.5cm]B) --([yshift=-0.5cm,xshift=-0.5cm]B)--([yshift=0.5cm,xshift=-0.5cm]C)--([yshift=0.5cm,xshift=0.5cm]C)--([yshift=-0.5cm,xshift=0.5cm]C)-- ([yshift=-0.5cm,xshift=-0.5cm]D)--([yshift=0.5cm,xshift=-0.5cm]D)--([yshift=0.5cm,xshift=0.5cm]D)--([yshift=-0.5cm,xshift=0.5cm]A)--([yshift=-0.5cm,xshift=-0.5cm]A)--([yshift=0.5cm,xshift=-0.5cm]A);
			
			\draw ([yshift=-0.3cm,xshift=0.3cm]A)-- ([yshift=-0.3cm,xshift=-0.3cm]B)--
			([yshift=0.3cm,xshift=-0.3cm]B) --([yshift=0.3cm,xshift=0.3cm]B)--([yshift=-0.3cm,xshift=0.3cm]C)--([yshift=-0.3cm,xshift=-0.3cm]C)--([yshift=0.3cm,xshift=-0.3cm]C)-- ([yshift=0.3cm,xshift=0.3cm]D)--([yshift=-0.3cm,xshift=0.3cm]D)--([yshift=-0.3cm,xshift=-0.3cm]D)--([yshift=0.3cm,xshift=-0.3cm]A)--([yshift=0.3cm,xshift=0.3cm]A)--([yshift=-0.3cm,xshift=0.3cm]A);
		\end{tikzpicture}
		
		\begin{center}
			\vspace{7pt}
			BỘ GIÁO DỤC VÀ ĐÀO TẠO
			
			\vspace{7pt}
			\textbf{TRƯỜNG ĐẠI HỌC TÂY BẮC}
		\end{center}
		
		\begin{center}
			\includegraphics[scale=0.5]{images/tb.jpg}
			
			\vspace{10pt}
			\fontsize{18pt}{17pt}\selectfont 
			\textbf{ ĐỒ ÁN TỐT NGHIỆP}
			
			\vspace{10pt}
			% \textbf{\footnotesize MÔN HỌC}
			\vspace{7pt}
			% \textbf{Môn học: Các vấn đề hiện đại của Công nghệ thông tin}
		\end{center}
		
		%		\begin{flushleft}
			%			\fontsize{14pt}{17pt}\selectfont
			%			\textbf{\textsl{ĐỀ TÀI}}
			%		\end{flushleft}
		
		\begin{center}
			\fontsize{18pt}{17pt}\selectfont
			\textbf{\textrm{XÂY DỰNG ỨNG DỤNG QUẢN LÝ KÝ TÚC XÁ TRƯỜNG ĐẠI HỌC TÂY BẮC}}
		\end{center}
		
		\vspace{75pt}
		\textbf{GV hướng dẫn: ThS. Phạm Quang Trung}
		
		\vspace{10pt}
		
		\textbf{Sinh viên thực hiện: Lành Văn Quyết}
		
		\vspace{10pt}
		\textbf{MSV: 2022A0871}
		
		\vspace{100pt}
		\begin{center}
			Sơn La - năm 2026
		\end{center}
	\end{titlepage}
	
	%===========================
	% Lời cảm ơn
	%===========================
	\newpage
	\thispagestyle{empty}
	\begin{center}
		\fontsize{14pt}{20pt}\selectfont
		\textbf{LỜI CẢM ƠN}
	\end{center}
	
	\vspace{20pt}
	
	Trong suốt quá trình nghiên cứu, phân tích và hoàn thiện báo cáo về hệ thống quản lý ký túc xá, em đã nhận được nhiều sự hỗ trợ quý báu từ thầy cô, bạn bè và những tài liệu chuyên môn liên quan đến phát triển phần mềm, quản lý cơ sở dữ liệu và xây dựng ứng dụng web hiện đại.
	Trước hết, em xin bày tỏ lòng biết ơn sâu sắc đến giảng viên hướng dẫn \textbf{ThS. Phạm Quang Trung} đã tận tình định hướng phương pháp tiếp cận, góp ý chuyên sâu về mặt học thuật và đồng hành cùng em trong suốt quá trình hoàn thiện nội dung khóa luận. Những nhận xét thẳng thắn và sát thực tế của thầy là cơ sở quan trọng giúp em điều chỉnh hướng đi và hoàn thiện hệ thống đúng với yêu cầu thực tiễn.
	Em xin chân thành cảm ơn Ban Giám hiệu \textbf{Trường Đại học Tây Bắc}, Ban chủ nhiệm \textbf{Khoa Công nghệ Thông tin} cùng toàn thể quý thầy cô đã truyền đạt kiến thức nền tảng về phân tích và thiết kế hệ thống, cơ sở dữ liệu, lập trình ứng dụng web và kiểm thử phần mềm trong suốt những năm học vừa qua. Đây chính là hành trang thiết yếu giúp em có thể tiếp cận bài toán quản lý ký túc xá một cách có hệ thống và khoa học.
	Em cũng xin gửi lời cảm ơn chân thành đến Ban Quản lý Ký túc xá \textbf{Trường Đại học Tây Bắc} đã tạo điều kiện để em tìm hiểu thực tế quy trình vận hành, từ đó xác định đúng các yêu cầu nghiệp vụ và xây dựng hệ thống phù hợp với nhu cầu sử dụng thực tế.
	Em xin cảm ơn gia đình và bạn bè đã luôn động viên, chia sẻ và tạo điều kiện thuận lợi trong suốt quá trình học tập và thực hiện đề tài. Sự ủng hộ đó là nguồn động lực giúp em duy trì sự tập trung và hoàn thành tốt nhiệm vụ được giao.
	Mặc dù đã cố gắng nghiên cứu và trình bày đầy đủ, báo cáo vẫn không tránh khỏi những thiếu sót nhất định do giới hạn về thời gian và kinh nghiệm thực tế. Em rất mong nhận được thêm các góp ý từ quý thầy cô và bạn đọc để nội dung được hoàn thiện hơn trong các lần cập nhật tiếp theo.
	
	\vspace{30pt}
	
	\begin{flushright}
		\textit{Sơn La, \today}\\[10pt]
		\textit{Sinh viên thực hiện}\\[30pt]
		\textbf{Lành Văn Quyết}
	\end{flushright}
	
	\newpage
	
	\newpage
	\begin{center}
		\textbf{NHẬN XÉT CỦA GIẢNG VIÊN HƯỚNG DẪN}
	\end{center}
	
	\vspace{1cm}
	
	\begin{enumerate}
		\item Hình thức (Bố cục, trình bày, lỗi, các mục, hình, bảng, công thức, phụ lục)
		
		\vspace{3cm}
		
		\item Nội dung (mục tiêu, phương pháp, kết quả, sao chép, các chương, tài liệu,..)
		
		\vspace{3cm}
		
		\item Kết luận
		
		\vspace{3cm}
	\end{enumerate}
	
	\vspace{1cm}
	
	\begin{flushright}
		\textit{Sơn La, Ngày ...... tháng ...... năm 202}\\[0.5cm]
		\textbf{Giảng viên hướng dẫn}\\[0.3cm]
		\textit{(Ký tên, ghi rõ họ tên)}
	\end{flushright}
	
	\newpage
	
	%	Nhận xét giảng viên
	\newpage
	\begin{center}
		\textbf{NHẬN XÉT CỦA GIẢNG VIÊN CHẤM}
	\end{center}
	
	\vspace{1cm}
	
	\begin{enumerate}
		\item Hình thức (Bố cục, trình bày, lỗi, các mục, hình, bảng, công thức, phụ lục)
		
		\vspace{3cm}
		
		\item Nội dung (mục tiêu, phương pháp, kết quả, sao chép, các chương, tài liệu,..)
		
		\vspace{3cm}
		
		\item Kết luận
		
		\vspace{3cm}
	\end{enumerate}
	
	\vspace{1cm}
	
	\begin{flushright}
		\textit{Sơn La, Ngày ...... tháng ...... năm 202}\\[0.5cm]
		\textbf{Giảng viên chấm}\\[0.3cm]
		\textit{(Ký tên, ghi rõ họ tên)}
	\end{flushright}
	\newpage
	
	%===========================
	% Mục lục
	%===========================
	\tableofcontents
	\thispagestyle{empty}
	\clearpage
	\newpage
	
	%===========================
	% Mở đầu
	%===========================
	\begin{abstract}
		\pagenumbering{arabic}
		\pagestyle{fancy}
		\fancyhf{}
		\fancyhead[L]{Khóa luận tốt nghiệp}
		\fancyhead[R]{\thepage}
		\renewcommand{\headrulewidth}{0.4pt}
		
		Trường Đại học Tây Bắc là cơ sở đào tạo đại học đa ngành trực thuộc Bộ Giáo dục và Đào tạo, với quy mô đào tạo hơn 5.000 sinh viên, học viên mỗi năm học. Trong đó, tỷ lệ sinh viên là người dân tộc thiểu số chiếm hơn 76\% và gần 500 lưu học sinh đến từ nước CHDCND Lào. Với đặc thù sinh viên phần lớn đến từ các tỉnh miền núi xa xôi, hệ thống ký túc xá của nhà trường đóng vai trò thiết yếu trong việc bảo đảm điều kiện học tập và sinh hoạt cho người học. Tuy nhiên, công tác quản lý nội trú hiện tại chủ yếu vẫn dựa trên phương thức thủ công hoặc các công cụ rời rạc, dẫn đến nhiều khó khăn trong việc theo dõi thông tin cư trú, quản lý hóa đơn, xử lý sự cố cơ sở vật chất và tổng hợp báo cáo vận hành.
		
		Đề tài này tập trung phân tích, thiết kế và xây dựng \textbf{Ứng dụng Quản lý Ký túc xá Trường Đại học Tây Bắc} -- một ứng dụng web hỗ trợ số hóa các nghiệp vụ cốt lõi trong công tác nội trú sinh viên. Hệ thống phục vụ ba vai trò chính là quản trị viên, nhân viên quản lý và sinh viên. Các chức năng chính bao gồm: xác thực và phân quyền người dùng, quản lý danh mục tòa nhà và phòng ở, tiếp nhận và xét duyệt đơn đăng ký nội trú, lập và theo dõi hóa đơn tiền phòng theo tháng, tiếp nhận yêu cầu bảo trì, gửi thông báo nội bộ và xem báo cáo thống kê vận hành. Phần backend được xây dựng bằng FastAPI, sử dụng PostgreSQL làm cơ sở dữ liệu; phần frontend được phát triển với React, TypeScript và Vite; toàn bộ hệ thống được đóng gói bằng Docker Compose để thuận lợi cho triển khai.
		
		\textbf{Nội dung, phạm vi của đề tài}
		\begin{itemize}
			\item \textbf{Nội dung:} Khảo sát bài toán quản lý nội trú sinh viên tại Đại học Tây Bắc, phân tích yêu cầu, thiết kế cơ sở dữ liệu và xây dựng hệ thống web quản lý phòng ở, đăng ký nội trú, hóa đơn, bảo trì và báo cáo thống kê.
			\item \textbf{Phạm vi:} Tập trung vào phiên bản web phục vụ quản trị viên, nhân viên và sinh viên; hóa đơn bao gồm tiền phòng, điện và nước; hệ thống hỗ trợ phân loại phòng theo đối tượng cư trú (giới tính, quốc tịch) và phân quyền theo vai trò. Phiên bản hiện tại chưa tích hợp thanh toán trực tuyến qua ngân hàng và chưa đồng bộ với hệ thống quản lý đào tạo của nhà trường.
		\end{itemize}
	\end{abstract}
	
	\clearpage
	\listoffigures
	\clearpage
	\listoftables
	
	
	\newpage
	\phantomsection
	%	\addcontentsline{toc}{section}{DANH MỤC TỪ VIẾT TẮT}
	\begin{center}
		\fontsize{14pt}{17pt}\selectfont
		{DANH MỤC TỪ VIẾT TẮT}
	\end{center}
	
	\vspace{10pt}
	
	\begin{tabular}{|c|p{11cm}|}
		\hline
		{Chữ viết tắt} & \multicolumn{1}{c|}{\textbf{Ý nghĩa}} \\
		\hline
		API  & Application Programming Interface, giao diện lập trình ứng dụng kết nối frontend và backend. \\
		\hline
		BE   & Backend, phần xử lý nghiệp vụ và truy xuất dữ liệu phía máy chủ. \\
		\hline
		CNTT & Công nghệ Thông tin. \\
		\hline
		CRUD & Create, Read, Update, Delete — bốn thao tác cơ bản với dữ liệu. \\
		\hline
		CSDL & Cơ sở dữ liệu. \\
		\hline
		ĐH   & Đại học. \\
		\hline
		FE   & Frontend, giao diện người dùng được xây dựng bằng React, TypeScript và Vite. \\
		\hline
		JWT  & JSON Web Token, chuẩn token dùng cho xác thực và phân quyền người dùng. \\
		\hline
		KTX  & Ký túc xá. \\
		\hline
		ORM  & Object Relational Mapping, cơ chế ánh xạ đối tượng với bảng trong cơ sở dữ liệu. \\
		\hline
		REST & Representational State Transfer, kiểu kiến trúc API được sử dụng trong hệ thống. \\
		\hline
		SV   & Sinh viên. \\
		\hline
		UUID & Universally Unique Identifier, mã định danh duy nhất toàn cục dùng cho các bản ghi. \\
		\hline
	\end{tabular}
	
	% ============================================================
	%  MỞ ĐẦU  (không đánh số, không dùng \section{} có số)
	% ============================================================
	
	\newpage
	\section*{MỞ ĐẦU}
	\addcontentsline{toc}{section}{MỞ ĐẦU}
	
	% -----------------------------------------------------------
	\subsection*{1. Lý do chọn đề tài}
	\addcontentsline{toc}{subsection}{1. Lý do chọn đề tài}
	
	Trong bối cảnh chuyển đổi số đang diễn ra mạnh mẽ tại các cơ sở giáo dục đại học Việt Nam, việc số hóa các quy trình quản lý hành chính -- trong đó có công tác quản lý ký túc xá -- ngày càng trở thành yêu cầu cấp thiết. Với quy mô đào tạo hơn 5.000 sinh viên và học viên mỗi năm học, trong đó tỷ lệ sinh viên là người dân tộc thiểu số chiếm hơn 76\% và gần 500 lưu học sinh đến từ nước Cộng hòa Dân chủ Nhân dân Lào \cite{utb_khaigiang2024}, Trường Đại học Tây Bắc có nhu cầu nội trú sinh viên rất lớn và mang nhiều đặc thù riêng biệt.
	
	Khu ký túc xá của nhà trường hiện có 8 tòa nhà 5 tầng với 473 phòng ở, đáp ứng khoảng 3.000 chỗ ở \cite{utb_ktx}. Tuy nhiên, công tác quản lý nội trú hiện tại chủ yếu vẫn dựa trên phương thức thủ công và các công cụ rời rạc, chưa có một hệ thống phần mềm thống nhất hỗ trợ toàn bộ chu trình từ đăng ký đến thanh lý nội trú. Thực trạng này dẫn đến nhiều khó khăn trong việc theo dõi thông tin cư trú, quản lý hóa đơn, xử lý sự cố cơ sở vật chất và tổng hợp báo cáo vận hành -- đặc biệt khi đối tượng cư trú đa dạng và số lượng người ở đông như tại Trường Đại học Tây Bắc.
	
	Xuất phát từ thực tiễn đó, đề tài \textbf{``Xây dựng hệ thống quản lý ký túc xá Trường Đại học Tây Bắc''} được lựa chọn nhằm phân tích, thiết kế và triển khai một ứng dụng web hỗ trợ số hóa các nghiệp vụ cốt lõi trong công tác nội trú sinh viên. Đây là bài toán có giá trị thực tiễn cao, đủ rõ về phạm vi và đủ phong phú về nghiệp vụ để xây dựng thành một hệ thống thông tin hoàn chỉnh, đồng thời cho phép vận dụng tổng hợp các kiến thức về phân tích yêu cầu, thiết kế cơ sở dữ liệu, xây dựng API RESTful, phát triển giao diện web và triển khai hệ thống trong môi trường container hóa.
	
	% -----------------------------------------------------------
	\subsection*{2. Mục tiêu nghiên cứu}
	\addcontentsline{toc}{subsection}{2. Mục tiêu nghiên cứu}
	
	Mục tiêu tổng quát của đề tài là xây dựng một hệ thống quản lý ký túc xá trên nền tảng web, hỗ trợ số hóa các nghiệp vụ cốt lõi trong công tác quản lý nội trú sinh viên. Hệ thống không chỉ đóng vai trò lưu trữ dữ liệu mà còn phải trở thành công cụ kiểm soát quy trình, giúp giảm thiểu thao tác thủ công, hạn chế sai sót nghiệp vụ và tạo môi trường tương tác minh bạch giữa sinh viên với bộ phận quản lý.
	
	Các mục tiêu cụ thể của đề tài bao gồm:
	
	\begin{itemize}
		\item Khảo sát, phân tích yêu cầu và xây dựng đặc tả hệ thống quản lý ký túc xá phù hợp với thực tiễn vận hành tại Trường Đại học Tây Bắc.
		\item Thiết kế cơ sở dữ liệu quan hệ và kiến trúc hệ thống theo mô hình nhiều lớp, đảm bảo tính nhất quán dữ liệu và khả năng mở rộng.
		\item Hiện thực hóa các chức năng cốt lõi bao gồm: xác thực và phân quyền người dùng; quản lý danh mục tòa nhà và phòng ở; tiếp nhận và xét duyệt đơn đăng ký nội trú; lập và theo dõi hóa đơn định kỳ; tiếp nhận và xử lý yêu cầu bảo trì; gửi thông báo nội bộ và xem báo cáo thống kê vận hành.
		\item Triển khai hệ thống trong môi trường container hóa sử dụng Docker Compose, đảm bảo tính ổn định và khả năng tái triển khai.
	\end{itemize}
	
	% -----------------------------------------------------------
	\subsection*{3. Đối tượng và phạm vi nghiên cứu}
	\addcontentsline{toc}{subsection}{3. Đối tượng và phạm vi nghiên cứu}
	
	\textbf{Đối tượng nghiên cứu} của đề tài là bài toán quản lý ký túc xá trong trường đại học, tập trung vào các quy trình nghiệp vụ phát sinh trong đời sống nội trú của sinh viên và công tác điều hành của bộ phận quản lý. Đối tượng nghiên cứu bao gồm ba nhóm chủ thể: sinh viên (người sử dụng dịch vụ nội trú), nhân viên quản lý (người vận hành quy trình) và quản trị viên hệ thống (người quản lý danh mục chung). Đề tài cũng nghiên cứu cách biểu diễn các quy tắc nghiệp vụ thành logic xử lý trong phần mềm, bao gồm điều kiện phê duyệt đơn nội trú, cơ chế cập nhật sức chứa phòng, điều kiện lập hóa đơn và phạm vi hợp lệ của yêu cầu bảo trì.
	
	\textbf{Phạm vi nghiên cứu} được xác định như sau:
	
	\begin{itemize}
		\item \textit{Phạm vi chức năng:} Đề tài tập trung vào các nghiệp vụ chính gồm xác thực người dùng, quản lý phòng ở, quản lý đăng ký nội trú, quản lý hóa đơn (tiền phòng, điện, nước), quản lý yêu cầu bảo trì, gửi thông báo và báo cáo thống kê theo tháng. Phiên bản hiện tại chưa tích hợp thanh toán trực tuyến qua cổng ngân hàng và chưa đồng bộ với hệ thống quản lý đào tạo của nhà trường.
		\item \textit{Phạm vi công nghệ:} Backend sử dụng FastAPI, SQLAlchemy, Alembic và PostgreSQL; frontend sử dụng React, Vite và TypeScript; môi trường triển khai sử dụng Docker Compose.
		\item \textit{Phạm vi người dùng:} Hệ thống phục vụ ba vai trò là quản trị viên, nhân viên quản lý và sinh viên, với cơ chế phân quyền theo vai trò được áp dụng xuyên suốt.
	\end{itemize}
	
	% -----------------------------------------------------------
	\subsection*{4. Phương pháp nghiên cứu}
	\addcontentsline{toc}{subsection}{4. Phương pháp nghiên cứu}
	
	Để thực hiện đề tài, tác giả kết hợp bốn nhóm phương pháp nghiên cứu chính, được triển khai song song thay vì tuần tự để đảm bảo tính gắn kết giữa lý thuyết và thực hành.
	
	\textbf{Phương pháp khảo sát nghiệp vụ} được áp dụng thông qua việc tổng hợp các tình huống quản lý thường gặp trong ký túc xá, phân tích nhu cầu thông tin của từng nhóm người dùng và đối chiếu với thực tế vận hành tại khu nội trú Trường Đại học Tây Bắc. Kết quả khảo sát làm cơ sở trực tiếp cho việc xác định yêu cầu chức năng và phi chức năng của hệ thống.
	
	\textbf{Phương pháp phân tích hệ thống} sử dụng kỹ thuật mô hình hóa tác nhân và ca sử dụng (use case) để bóc tách bài toán thành các luồng nghiệp vụ có thể kiểm soát. Song song với đó, tư duy phân tích trạng thái (state analysis) được vận dụng để theo dõi sự chuyển đổi của các đối tượng nghiệp vụ quan trọng như đơn đăng ký, phòng ở, hóa đơn và yêu cầu bảo trì -- nhằm tránh thiết kế hệ thống chỉ dựa trên giao diện mà thiếu đi các ràng buộc xử lý bên trong.
	
	\textbf{Phương pháp thiết kế hệ thống} được thực hiện theo kiến trúc nhiều lớp, tách biệt giao diện, xử lý nghiệp vụ và dữ liệu. Ngôn ngữ mô hình hóa thống nhất (UML) được sử dụng để mô tả hành vi hệ thống thông qua biểu đồ ca sử dụng, biểu đồ hoạt động và biểu đồ tuần tự -- phù hợp với mục tiêu làm rõ quy trình nghiệp vụ và tương tác giữa các thành phần.
	
	\textbf{Phương pháp lập trình và thử nghiệm} được tiến hành bám sát các quy tắc nghiệp vụ đã xác định trước đó. Việc thử nghiệm được thực hiện thông qua các kịch bản sử dụng điển hình, nhằm kiểm tra tính đúng đắn của dữ liệu đầu vào, tính hợp lệ của trạng thái nghiệp vụ và sự phối hợp chính xác giữa frontend và backend.
	
	% -----------------------------------------------------------
	\subsection*{5. Cấu trúc báo cáo}
	\addcontentsline{toc}{subsection}{5. Cấu trúc báo cáo}
	
	Ngoài phần Mở đầu, Kết luận và các danh mục phụ lục, báo cáo được tổ chức thành bốn chương chính:
	
	\textbf{Chương 1 -- Khảo sát và Đặc tả hệ thống:} Trình bày bối cảnh và lý do chọn đề tài; kết quả khảo sát thực tế tại khu ký túc xá Trường Đại học Tây Bắc; phân tích yêu cầu chức năng và phi chức năng; đặc tả hệ thống thông qua các biểu đồ ca sử dụng và mô tả luồng nghiệp vụ.
	
	\textbf{Chương 2 -- Phân tích và Thiết kế hệ thống:} Trình bày kiến trúc tổng thể của hệ thống; thiết kế cơ sở dữ liệu; thiết kế các luồng xử lý nghiệp vụ thông qua biểu đồ hoạt động và biểu đồ tuần tự; thiết kế giao diện người dùng.
	
	\textbf{Chương 3 -- Cài đặt và Triển khai hệ thống:} Mô tả môi trường phát triển và công nghệ sử dụng; trình bày chi tiết cài đặt các chức năng cốt lõi ở cả backend và frontend; hướng dẫn triển khai hệ thống với Docker Compose.
	
	\textbf{Chương 4 -- Kiểm thử và Đánh giá:} Trình bày kế hoạch kiểm thử, các kịch bản thử nghiệm điển hình và kết quả kiểm thử; đánh giá mức độ đáp ứng yêu cầu và định hướng phát triển tiếp theo của hệ thống.
	
	% ============================================================
	%  CHƯƠNG 1: KHẢO SÁT VÀ ĐẶC TẢ HỆ THỐNG
	% ============================================================
	
	\newpage
	\section{KHẢO SÁT VÀ ĐẶC TẢ HỆ THỐNG}
	
	% -----------------------------------------------------------
	\subsection{Bối cảnh và lý do chọn đề tài}
	
	Trường Đại học Tây Bắc được thành lập theo Quyết định số 39/2001/QĐ-TTg ngày 23 tháng 3 năm 2001 của Thủ tướng Chính phủ, trên cơ sở nâng cấp Trường Cao đẳng Sư phạm Tây Bắc \cite{utb_history}. Đây là trường đại học đa ngành duy nhất của khu vực Tây Bắc, đặt trụ sở tại thành phố Sơn La, với sứ mệnh đào tạo nguồn nhân lực chất lượng cao phục vụ phát triển kinh tế -- xã hội vùng Tây Bắc. Quy mô đào tạo hằng năm đạt hơn 5.000 học viên và sinh viên, trong đó tỷ lệ sinh viên là người dân tộc thiểu số chiếm hơn 76\% \cite{utb_khaigiang2024}. Đặc thù sinh viên phần lớn đến từ các huyện miền núi xa xôi thuộc các tỉnh Sơn La, Điện Biên, Lai Châu và từ nước bạn Lào tạo ra nhu cầu nội trú rất lớn và mang tính đặc thù cao.
	
	Khu ký túc xá của nhà trường hiện bao gồm 8 tòa nhà 5 tầng với 473 phòng ở và 8 phòng sinh hoạt chung, đáp ứng khoảng 3.000 chỗ ở với mức phí chỉ từ 50.000 đồng/người/tháng \cite{utb_ktx}. Mỗi phòng bố trí từ 6 đến 8 sinh viên, được trang bị đầy đủ tủ sắt, giường ngủ, bàn học, quạt trần và công trình vệ sinh khép kín \cite{utb_phong_ktu}. Với gần 500 lưu học sinh Lào sinh sống trong khu nội trú cùng hàng nghìn sinh viên Việt Nam thuộc nhiều dân tộc khác nhau, công tác quản lý ký túc xá đặt ra những yêu cầu cao về phân loại đối tượng cư trú, theo dõi thông tin đa dạng và kiểm soát trật tự nội trú. Đơn vị trực tiếp phụ trách công tác này là Ban Quản lý khu nội trú thuộc Phòng Quản trị -- Thiết bị, đặt tại ký túc xá K2 của trường \cite{utb_phongsban}.
	
	Tuy nhiên, khảo sát thực tiễn cho thấy công tác quản lý nội trú tại nhà trường hiện nay vẫn chủ yếu dựa vào bảng tính điện tử, hồ sơ giấy và trao đổi trực tiếp giữa sinh viên với bộ phận quản lý, chưa có một phần mềm thống nhất hỗ trợ toàn bộ chu trình từ đăng ký đến thanh lý nội trú. Khi dữ liệu về phòng ở, hóa đơn và phản ánh bảo trì nằm ở các kênh riêng rẽ, bộ phận quản lý gặp khó khăn trong việc kiểm soát tình trạng vận hành chung. Đồng thời, sinh viên cũng thiếu kênh chủ động để theo dõi đơn đăng ký, kiểm tra khoản phải nộp hoặc nắm bắt tình trạng xử lý các yêu cầu đã gửi.
	
	Chính vì vậy, việc xây dựng một hệ thống quản lý ký túc xá có tính tập trung, nhất quán và minh bạch là nhu cầu thực sự cần thiết, đồng thời mang lại giá trị ứng dụng cao đối với Trường Đại học Tây Bắc trong giai đoạn hiện nay.
	
	% -----------------------------------------------------------
	\subsection{Khảo sát thực tế}
	
	\subsubsection{Kết quả khảo sát}
	
	Khảo sát nghiệp vụ tại khu ký túc xá Trường Đại học Tây Bắc cho thấy hoạt động quản lý nội trú xoay quanh một chuỗi quy trình lặp lại theo chu kỳ, có thể phân thành bốn nhóm nghiệp vụ chính.
	
	\textbf{Quản lý hạ tầng lưu trú:} Bộ phận phụ trách cần nắm được tình trạng hoạt động của từng tòa nhà, số lượng phòng theo từng khu, đặc tính của từng phòng (nhóm đối tượng được phép ở, sức chứa, trạng thái sử dụng) và số chỗ còn trống tại từng thời điểm. Việc phân loại phòng theo đối tượng (nam, nữ, lưu học sinh Lào) là yêu cầu đặc thù riêng của trường, không thể bỏ qua trong quá trình thiết kế hệ thống.
	
	\textbf{Tiếp nhận và quản lý cư trú:} Quy trình tiếp nhận sinh viên nội trú bao gồm các bước xem phòng phù hợp, nộp đơn đăng ký, xét duyệt, bố trí vào ở và cập nhật tình trạng cư trú. Mỗi bước đều phát sinh điều kiện kiểm tra và hành động cập nhật dữ liệu liên quan, tạo thành một luồng nghiệp vụ có trật tự chặt chẽ. Khi sinh viên kết thúc thời gian ở, quy trình thanh lý cư trú cũng cần được ghi nhận để cập nhật lại sức chứa phòng và trạng thái đăng ký.
	
	\textbf{Quản lý hóa đơn và thanh toán:} Hóa đơn được lập định kỳ theo tháng, bao gồm tiền phòng, tiền điện và tiền nước. Bộ phận quản lý cần theo dõi tình trạng thanh toán của từng sinh viên, xác nhận khi đã nhận tiền và phát hiện các trường hợp chưa thanh toán để có biện pháp nhắc nhở phù hợp.
	
	\textbf{Bảo trì và thông báo:} Sinh viên cần có kênh gửi phản ánh về hỏng hóc hoặc sự cố trong phòng ở. Bộ phận quản lý tiếp nhận, phân loại và xử lý từng yêu cầu theo mức độ ưu tiên. Song song với đó, ban quản lý cũng có nhu cầu gửi thông báo chung (lịch đóng phí, quy định mới, thông tin sự kiện) đến đúng nhóm đối tượng.
	
	Hiện tại, toàn bộ các công việc trên chủ yếu được thực hiện thủ công hoặc qua bảng tính, chưa có một phần mềm thống nhất hỗ trợ toàn bộ chu trình từ đăng ký đến thanh lý nội trú \cite{utb_phongsban}. Điều này dẫn đến tình trạng dữ liệu phân tán, khó tổng hợp báo cáo và chậm trễ trong phản hồi các yêu cầu phát sinh của sinh viên.
	
	\subsubsection{Phân tích kết quả khảo sát}
	
	Từ kết quả khảo sát, có thể nhận thấy bài toán quản lý ký túc xá ẩn chứa sự phức tạp đáng kể khi nhìn từ góc độ luồng dữ liệu và trạng thái nghiệp vụ. Một thao tác tưởng như đơn giản -- chẳng hạn phê duyệt cho sinh viên vào ở -- thực tế kéo theo nhiều điều kiện kiểm tra đồng thời: tính phù hợp của loại phòng với đối tượng sinh viên, trạng thái hiện tại của phòng (còn chỗ hay đã đầy, đang bảo trì hay đang hoạt động bình thường), và tình trạng đăng ký hiện tại của sinh viên (đã có đăng ký đang có hiệu lực chưa). Tương tự, việc lập hóa đơn cũng phụ thuộc vào quan hệ hợp lệ giữa sinh viên với phòng ở và kỳ thanh toán, chứ không thể tạo tùy ý nếu chưa có cơ sở cư trú hợp lệ.
	
	Phân tích này dẫn đến một yêu cầu thiết kế quan trọng: hệ thống không thể chỉ là giao diện nhập liệu đơn thuần, mà phải được xây dựng theo hướng \textit{kiểm soát quy trình}, trong đó tầng nghiệp vụ (service layer) đóng vai trò trung tâm để kiểm tra điều kiện, thực thi ràng buộc và cập nhật trạng thái một cách nhất quán. Đây là điểm mấu chốt để hệ thống vừa đảm bảo tính đúng đắn của dữ liệu, vừa thực sự hữu ích trong vận hành thực tế.
	
	Bên cạnh đó, sự đa dạng của đối tượng cư trú tại Trường Đại học Tây Bắc đặt ra yêu cầu hệ thống phải xử lý được các trường hợp phòng dành riêng theo nhóm, không thể áp dụng một mô hình quản lý phòng đồng nhất. Đây là điểm khác biệt quan trọng so với bài toán ký túc xá tại các cơ sở đào tạo thông thường.
	
	% -----------------------------------------------------------
	\subsection{Đặc tả hệ thống}
	
	\subsubsection{Tổng quan kiến trúc}
	
	Hệ thống quản lý ký túc xá được đặc tả như một ứng dụng web nhiều người dùng, hoạt động theo mô hình client-server. Phía client, sinh viên và cán bộ quản lý thao tác thông qua giao diện web được xây dựng bằng React và TypeScript. Phía server, backend tiếp nhận yêu cầu thông qua API RESTful, thực hiện xác thực danh tính, kiểm tra quyền truy cập, áp dụng các quy tắc nghiệp vụ và truy xuất dữ liệu từ cơ sở dữ liệu quan hệ PostgreSQL. Toàn bộ hệ thống được đóng gói bằng Docker Compose để đảm bảo tính nhất quán giữa các môi trường triển khai.
	
	\subsubsection{Các tác nhân và vai trò}
	
	Hệ thống xác định ba nhóm tác nhân với phạm vi thao tác được phân quyền rõ ràng:
	
	\begin{itemize}
		\item \textbf{Quản trị viên (Admin):} Quản lý danh mục chung của hệ thống bao gồm thông tin người dùng, tòa nhà và phòng ở. Quản trị viên có quyền truy cập rộng nhất, bao gồm cả việc theo dõi báo cáo tổng hợp và giám sát vận hành toàn hệ thống.
		\item \textbf{Nhân viên quản lý (Staff):} Là người trực tiếp vận hành các quy trình nghiệp vụ hằng ngày: xét duyệt đơn đăng ký nội trú, lập và xác nhận hóa đơn, xử lý yêu cầu bảo trì và gửi thông báo đến sinh viên.
		\item \textbf{Sinh viên (Student):} Là người dùng cuối sử dụng hệ thống để tra cứu thông tin phòng còn chỗ, nộp đơn đăng ký, theo dõi trạng thái xét duyệt, kiểm tra hóa đơn và gửi yêu cầu bảo trì.
	\end{itemize}
	
	\subsubsection{Yêu cầu chức năng}
	
	Hệ thống bao gồm các nhóm chức năng chính sau:
	
	\begin{enumerate}
		\item \textbf{Xác thực và phân quyền:} Đăng nhập bằng tên tài khoản và mật khẩu; xác thực theo cơ chế JWT (JSON Web Token); phân quyền theo vai trò được áp dụng xuyên suốt tất cả các API.
		\item \textbf{Quản lý tòa nhà và phòng ở:} Thêm, sửa, xóa thông tin tòa nhà; quản lý danh sách phòng theo từng tòa, bao gồm sức chứa, loại đối tượng được phép ở và trạng thái phòng (đang sử dụng, đang bảo trì, còn trống).
		\item \textbf{Đăng ký và quản lý nội trú:} Sinh viên nộp đơn đăng ký phòng; nhân viên xét duyệt; hệ thống tự động cập nhật số người ở thực tế và trạng thái phòng sau mỗi quyết định; xử lý thanh lý khi sinh viên rời khỏi ký túc xá.
		\item \textbf{Quản lý hóa đơn:} Lập hóa đơn định kỳ theo tháng (tiền phòng, điện, nước); theo dõi trạng thái thanh toán; xác nhận khi nhận tiền; thống kê công nợ theo kỳ.
		\item \textbf{Quản lý bảo trì:} Tiếp nhận yêu cầu bảo trì từ sinh viên đang cư trú hợp lệ; cập nhật trạng thái xử lý; theo dõi lịch sử bảo trì theo phòng.
		\item \textbf{Thông báo:} Nhân viên và quản trị viên gửi thông báo đến cá nhân hoặc nhóm sinh viên; sinh viên xem và đánh dấu đã đọc.
		\item \textbf{Báo cáo thống kê:} Tổng hợp số liệu theo tháng về tỷ lệ lấp đầy từng khu nhà, tình hình thu phí và thống kê yêu cầu bảo trì.
	\end{enumerate}
	
	\subsubsection{Yêu cầu phi chức năng}
	
	\begin{itemize}
		\item \textbf{Bảo mật:} Mật khẩu được mã hóa bằng thuật toán bcrypt; token JWT có thời hạn sử dụng giới hạn; mọi endpoint đều yêu cầu xác thực và kiểm tra quyền trước khi xử lý.
		\item \textbf{Tính nhất quán dữ liệu:} Các thao tác liên quan đến trạng thái nghiệp vụ (duyệt đăng ký, lập hóa đơn, xử lý bảo trì) được thực thi trong phạm vi giao dịch cơ sở dữ liệu để đảm bảo toàn vẹn dữ liệu.
		\item \textbf{Khả năng mở rộng:} Kiến trúc nhiều lớp (router -- service -- repository) cho phép bổ sung tính năng mới mà không ảnh hưởng đến các thành phần hiện có.
		\item \textbf{Khả năng triển khai:} Hệ thống được đóng gói bằng Docker Compose, có thể triển khai nhất quán trên các môi trường khác nhau mà không cần cấu hình thủ công phức tạp.
		\item \textbf{Giao diện tiếng Việt:} Toàn bộ giao diện người dùng sử dụng tiếng Việt, phù hợp với ngữ cảnh sử dụng thực tế tại Trường Đại học Tây Bắc.
	\end{itemize}
	
	\subsubsection{Ràng buộc nghiệp vụ}
	
	Một số ràng buộc nghiệp vụ trọng tâm được đặc tả như sau:
	
	\begin{itemize}
		\item Sinh viên chỉ được đăng ký vào phòng phù hợp với nhóm đối tượng của mình (nam/nữ/lưu học sinh Lào).
		\item Một sinh viên không được có đồng thời nhiều hơn một đơn đăng ký đang có hiệu lực.
		\item Phòng đã đạt sức chứa tối đa hoặc đang trong trạng thái bảo trì không được tiếp nhận đăng ký mới.
		\item Hóa đơn không được tạo trùng cho cùng một sinh viên trong cùng một kỳ thanh toán.
		\item Yêu cầu bảo trì chỉ được gửi cho phòng mà sinh viên đang cư trú hợp lệ tại thời điểm gửi.
		\item Sau khi đăng ký được duyệt, số người ở thực tế trong phòng được cập nhật tự động; sau khi thanh lý, số người ở giảm tương ứng và trạng thái phòng được rà soát lại.
	\end{itemize}
	
	% Ghi chú: Sơ đồ ca sử dụng tổng quát và các biểu đồ chi tiết mô tả các luồng nghiệp vụ trên được trình bày trong Chương~2.
	
	\newpage
	\section*{CHƯƠNG 2: PHÂN TÍCH HỆ THỐNG}
	\addcontentsline{toc}{section}{CHƯƠNG 2: PHÂN TÍCH HỆ THỐNG}
	
	\subsection*{1. Mô hình hóa hệ thống}
	\addcontentsline{toc}{subsection}{1. Mô hình hóa hệ thống}
	
	% -----------------------------------------------------------
	\subsubsection*{1.1 Ca sử dụng tổng quát}
	\addcontentsline{toc}{subsubsection}{1.1 Ca sử dụng tổng quát}
	
	%	\begin{figure}[H]
		%		\centering
		%		\includegraphics[width=1\linewidth]{images/uc_uc00}
		%		\caption{Biểu đồ ca sử dụng tổng thể --- Hệ thống Quản lý Ký túc xá}
		%		\label{fig:ucuc00}
		%	\end{figure}
	\begin{figure}[H]
		\centering
		\includegraphics[width=0.8\linewidth]{images/uc_uc00}
		\caption{Biểu đồ ca sử dụng tổng quát}
		\label{fig:ucuc00}
	\end{figure}
	
	
	Biểu đồ tổng thể mô tả toàn bộ các ca sử dụng của hệ thống và mối quan hệ giữa ba nhóm tác nhân chính: Sinh viên, Nhân viên quản lý và Quản trị viên. Sinh viên thực hiện các thao tác liên quan đến việc sử dụng dịch vụ nội trú như đăng nhập, đăng ký phòng, thanh toán hóa đơn và gửi yêu cầu bảo trì. Nhân viên quản lý phụ trách các nghiệp vụ vận hành hằng ngày gồm duyệt đơn đăng ký, xử lý trả phòng, lập hóa đơn, xử lý bảo trì và gửi thông báo. Quản trị viên có quyền truy cập rộng nhất, bao gồm cấp tài khoản, quản lý tòa nhà và phòng ở, gửi thông báo và xem báo cáo thống kê vận hành.
	
	% -----------------------------------------------------------
	\subsubsection*{1.2 Ca sử dụng cấp tài khoản}
	\addcontentsline{toc}{subsubsection}{1.2 Ca sử dụng cấp tài khoản}
	
	\paragraph{Biểu đồ ca sử dụng}
	\noindent
	
	\begin{figure}[H]
		\centering
		\includegraphics[width=0.85\linewidth]{images/uc_uc01}
		\caption{Biểu đồ ca sử dụng cấp tài khoản}
		\label{fig:ucuc01}
	\end{figure}
	
	
	\begin{center}
		\begin{tabular}{|p{4cm}|p{10cm}|}
			\hline
			\textbf{Tên ca sử dụng} & \textbf{Cấp tài khoản người dùng} \\
			\hline
			Tác nhân chính & Quản trị viên \\
			\hline
			Dòng sự kiện chính &
			\begin{minipage}[t]{10cm}
				\vspace{2pt}
				\begin{enumerate}
					\item Quản trị viên truy cập chức năng Cấp tài khoản.
					\item Nhập thông tin: họ tên, email, mã số sinh viên/nhân viên, giới tính, quốc tịch và vai trò.
					\item Hệ thống kiểm tra tính hợp lệ và phát hiện trùng lặp.
					\item Tạo tài khoản với mật khẩu mặc định và gán vai trò tương ứng.
					\item Hiển thị xác nhận tạo tài khoản thành công.
				\end{enumerate}
				\vspace{2pt}
			\end{minipage} \\
			\hline
			Dòng sự kiện phụ & Nếu email hoặc mã số đã tồn tại trong hệ thống, hệ thống trả về lỗi 409 Conflict và yêu cầu nhập lại. Nếu dữ liệu không hợp lệ (thiếu trường bắt buộc, sai định dạng), hệ thống hiển thị thông báo lỗi xác thực tương ứng. \\
			\hline
			Điều kiện tiên quyết & Người thực hiện đã đăng nhập với vai trò Quản trị viên. \\
			\hline
			Điều kiện sau & Tài khoản mới được tạo thành công và có thể đăng nhập vào hệ thống với vai trò được gán. \\
			\hline
		\end{tabular}
	\end{center}
	
	\paragraph{Biểu đồ hoạt động}
	\noindent
	
	\begin{figure}[H]
		\centering
		\includegraphics[width=0.55\linewidth]{images/activity_uc01}
		\caption{Biểu đồ hoạt động cấp tài khoản}
		\label{fig:activityuc01}
	\end{figure}
	
	
	Quản trị viên điền thông tin người dùng mới; hệ thống kiểm tra tính hợp lệ của dữ liệu đầu vào và tra soát trùng lặp email hoặc mã số. Nếu hợp lệ, hệ thống tạo tài khoản với mật khẩu băm mặc định và gán vai trò; ngược lại trả thông báo lỗi tương ứng (lỗi xác thực hoặc 409 Conflict).
	
	\paragraph{Biểu đồ tuần tự}
	\noindent
	
	\begin{figure}[H]
		\centering
		\includegraphics[width=0.9\linewidth]{images/sequence_uc01}
		\caption{Biểu đồ tuần tự cấp tài khoản}
		\label{fig:sequenceuc01}
	\end{figure}
	
	
	Quản trị viên gửi \texttt{POST /api/v1/admin/users} kèm thông tin người dùng; Backend kiểm tra trùng lặp email và mã số trong Database, nếu hợp lệ thì băm mật khẩu mặc định, lưu bản ghi mới vào bảng \texttt{users} và trả về 201 Created với thông tin tài khoản đã tạo.
	
	% -----------------------------------------------------------
	\subsubsection*{1.3 Ca sử dụng đăng nhập}
	\addcontentsline{toc}{subsubsection}{1.3 Ca sử dụng đăng nhập}
	
	\paragraph{Biểu đồ ca sử dụng}
	\noindent
	
	\begin{figure}[H]
		\centering
		\includegraphics[width=0.85\linewidth]{images/uc_uc02}
		\caption{Biểu đồ ca sử dụng đăng nhập}
		\label{fig:ucuc02}
	\end{figure}
	
	
	\begin{center}
		\begin{tabular}{|p{4cm}|p{10cm}|}
			\hline
			\textbf{Tên ca sử dụng} & \textbf{Đăng nhập hệ thống} \\
			\hline
			Tác nhân chính & Người dùng đã được cấp tài khoản (Quản trị viên, Nhân viên hoặc Sinh viên) \\
			\hline
			Dòng sự kiện chính &
			\begin{minipage}[t]{10cm}
				\vspace{2pt}
				\begin{enumerate}
					\item Người dùng nhập email và mật khẩu trên màn hình đăng nhập.
					\item Hệ thống xác thực thông tin đăng nhập và kiểm tra trạng thái tài khoản.
					\item Nếu hợp lệ, hệ thống phát hành JWT Access Token và Refresh Token lưu trong HttpOnly cookie.
					\item Hệ thống điều hướng đến Dashboard tương ứng với vai trò người dùng.
				\end{enumerate}
				\vspace{2pt}
			\end{minipage} \\
			\hline
			Dòng sự kiện phụ & Nếu email không tồn tại hoặc mật khẩu sai, hệ thống trả thông báo lỗi 401 Unauthorized. Khi Access Token hết hạn, hệ thống tự động làm mới bằng Refresh Token thông qua cơ chế \textit{token rotation}. \\
			\hline
			Điều kiện tiên quyết & Người dùng đã có tài khoản được Quản trị viên tạo trong hệ thống. \\
			\hline
			Điều kiện sau & Người dùng đăng nhập thành công và có quyền truy cập các chức năng tương ứng với vai trò. \\
			\hline
		\end{tabular}
	\end{center}
	
	\paragraph{Biểu đồ hoạt động}
	\noindent
	\begin{figure}[H]
		\centering
		\includegraphics[width=0.45\linewidth]{images/activity_uc02}
		\caption{Biểu đồ hoạt động đăng nhập}
		\label{fig:activityuc02}
	\end{figure}
	
	
	Người dùng nhập email và mật khẩu; hệ thống kiểm tra tính hợp lệ của thông tin xác thực, xác minh email tồn tại và đối chiếu mật khẩu băm. Nếu hợp lệ, hệ thống tạo cặp JWT Access Token và Refresh Token lưu trong HttpOnly cookie rồi chuyển hướng theo vai trò; ngược lại trả thông báo lỗi 401.
	
	\paragraph{Biểu đồ tuần tự}
	\noindent
	\begin{figure}[H]
		\centering
		\includegraphics[width=0.9\linewidth]{images/sequence_uc02}
		\caption{Biểu đồ tuần tự đăng nhập}
		\label{fig:sequenceuc02}
	\end{figure}
	
	Frontend gửi \texttt{POST /api/v1/auth/login} với email và mật khẩu đến Auth API; Backend tra cứu tài khoản trong Database theo email, xác thực mật khẩu băm bằng bcrypt, sau đó phát hành cặp JWT Token trong HttpOnly cookie và điều hướng Frontend đến Dashboard phù hợp với vai trò.
	
	% -----------------------------------------------------------
	\subsubsection*{1.4 Ca sử dụng uản lý tòa nhà và phòng ở}
	\addcontentsline{toc}{subsubsection}{1.4 Ca sử dụng quản lý tòa nhà và phòng ở}
	
	\paragraph{Biểu đồ ca sử dụng}
	\noindent	
	\begin{figure}[H]
		\centering
		\includegraphics[width=0.85\linewidth]{images/uc_uc03}
		\caption{Biểu đồ ca sử dụng quản lý tòa nhà và phòng ở}
		\label{fig:ucuc03}
	\end{figure}
	
	
	\begin{center}
		\begin{tabular}{|p{4cm}|p{10cm}|}
			\hline
			\textbf{Tên ca sử dụng} & \textbf{Quản lý tòa nhà và phòng ở} \\
			\hline
			Tác nhân chính & Quản trị viên \\
			\hline
			Dòng sự kiện chính &
			\begin{minipage}[t]{10cm}
				\vspace{2pt}
				\begin{enumerate}
					\item Quản trị viên chọn chức năng quản lý tòa nhà hoặc phòng ở.
					\item Thực hiện thao tác thêm mới, chỉnh sửa hoặc xóa tòa nhà/phòng.
					\item Hệ thống kiểm tra ràng buộc nghiệp vụ (không xóa khi còn sinh viên đang ở, không để sức chứa âm, v.v.).
					\item Cập nhật thông tin và phản hồi kết quả cho người dùng.
				\end{enumerate}
				\vspace{2pt}
			\end{minipage} \\
			\hline
			Dòng sự kiện phụ & Khi xóa tòa nhà còn phòng hoạt động hoặc xóa phòng còn sinh viên đang ở, hệ thống trả lỗi 400/409 và từ chối thao tác. Cập nhật trạng thái phòng (AVAILABLE, FULL, MAINTENANCE) được thực hiện tự động hoặc thủ công tùy tình huống. \\
			\hline
			Điều kiện tiên quyết & Người thực hiện đã đăng nhập với vai trò Quản trị viên. \\
			\hline
			Điều kiện sau & Danh mục tòa nhà và phòng ở được cập nhật nhất quán trong cơ sở dữ liệu. \\
			\hline
		\end{tabular}
	\end{center}
	
	\paragraph{Biểu đồ hoạt động}
	\noindent
	
	\begin{figure}[H]
		\centering
		\includegraphics[width=0.55\linewidth]{images/activity_uc03}
		\caption{Biểu đồ hoạt động quản lý tòa nhà và phòng ở}
		\label{fig:activityuc03}
	\end{figure}
	
	
	Quản trị viên chọn đối tượng quản lý (tòa nhà hoặc phòng) và loại thao tác (thêm/sửa/xóa); hệ thống xử lý song song hai luồng: quản lý tòa nhà và quản lý phòng, mỗi luồng đều kiểm tra ràng buộc trước khi thực thi. Kết quả thao tác được phản hồi và danh sách được cập nhật.
	
	\paragraph{Biểu đồ tuần tự}
	\noindent
	\begin{figure}[H]
		\centering
		\includegraphics[width=1\linewidth]{images/sequence_uc03}
		\caption{Biểu đồ tuần tự quản lý tòa nhà và phòng ở}
		\label{fig:sequenceuc03}
	\end{figure}
	
	
	Quản trị viên gửi yêu cầu \texttt{POST/PUT/DELETE} đến \texttt{/api/v1/buildings} hoặc \texttt{/rooms}; Backend kiểm tra quyền ADMIN, thực hiện thao tác CRUD trên Database. Nếu vi phạm ràng buộc, trả về 400/409 kèm thông báo lỗi; nếu thành công, trả về 200/201 với dữ liệu cập nhật.
	
	% -----------------------------------------------------------
	\subsubsection*{1.5 Ca sử dụng đăng ký phòng}
	\addcontentsline{toc}{subsubsection}{1.5 Ca sử dụng đăng ký phòng}
	
	\paragraph{Biểu đồ ca sử dụng}
	\noindent
	
	\begin{figure}[H]
		\centering
		\includegraphics[width=1\linewidth]{images/uc_uc05}
		\caption{Biểu đồ ca sử dụng đăng ký phòng}
		\label{fig:ucuc05}
	\end{figure}
	
	
	\begin{center}
		\begin{tabular}{|p{4cm}|p{10cm}|}
			\hline
			\textbf{Tên ca sử dụng} & \textbf{Đăng ký phòng ở ký túc xá} \\
			\hline
			Tác nhân chính & Sinh viên \\
			\hline
			Dòng sự kiện chính &
			\begin{minipage}[t]{10cm}
				\vspace{2pt}
				\begin{enumerate}
					\item Sinh viên xem danh sách phòng còn chỗ trống (trạng thái AVAILABLE).
					\item Hệ thống lọc và hiển thị các phòng phù hợp với giới tính và quốc tịch của sinh viên.
					\item Sinh viên chọn phòng và xác nhận đăng ký với ngày bắt đầu và ngày kết thúc dự kiến.
					\item Hệ thống tạo đơn đăng ký với trạng thái PENDING và thông báo cho sinh viên.
				\end{enumerate}
				\vspace{2pt}
			\end{minipage} \\
			\hline
			Dòng sự kiện phụ & Nếu sinh viên đã có đăng ký đang hoạt động (APPROVED), hệ thống từ chối và trả lỗi 400. Nếu phòng không phù hợp với loại đối tượng (sai giới tính/quốc tịch) hoặc phòng đã đầy, hệ thống cũng từ chối với thông báo cụ thể. \\
			\hline
			Điều kiện tiên quyết & Sinh viên đã đăng nhập và chưa có đăng ký đang có hiệu lực. \\
			\hline
			Điều kiện sau & Đơn đăng ký được tạo với trạng thái PENDING, chờ nhân viên xét duyệt. \\
			\hline
		\end{tabular}
	\end{center}
	
	\paragraph{Biểu đồ hoạt động}
	\noindent
	
	\begin{figure}[H]
		\centering
		\includegraphics[width=0.6\linewidth]{images/activity_uc05}
		\caption{Biểu đồ hoạt động đăng ký phòng}
		\label{fig:activityuc05}
	\end{figure}
	
	
	Sinh viên xem danh sách phòng AVAILABLE và chọn phòng; hệ thống kiểm tra tuần tự ba điều kiện: sinh viên chưa có đăng ký đang hoạt động, giới tính và quốc tịch phù hợp loại phòng, và phòng còn chỗ trống. Nếu vượt qua cả ba điều kiện, đơn đăng ký được tạo với trạng thái PENDING.
	
	\paragraph{Biểu đồ tuần tự}
	\noindent
	
	%	\begin{figure}[H]
		%		\centering
		%		\includegraphics[width=0.9\linewidth]{images/sequence_uc05}
		%		\caption{Biểu đồ tuần tự UC05 --- Đăng ký phòng}
		%		\label{fig:sequenceuc05}
		%	\end{figure}
	
	Sinh viên gửi \texttt{GET /api/v1/rooms?status=AVAILABLE} để lấy danh sách phòng, sau đó gửi \texttt{POST /api/v1/registrations} để đăng ký; Backend kiểm tra điều kiện sinh viên và phòng trong Database, nếu hợp lệ thì \texttt{INSERT} bản ghi đăng ký với trạng thái PENDING và trả về 201 Created.
	
	\begin{figure}[H]
		\centering
		\includegraphics[width=0.9\linewidth]{images/sequence_uc05}
		\caption{Biểu đồ tuần tự đăng ký phòng}
		\label{fig:sequenceuc05}
	\end{figure}
	
	
	% -----------------------------------------------------------
	\subsubsection*{1.6 Ca sử dụng duyệt đơn đăng ký}
	\addcontentsline{toc}{subsubsection}{1.6 Ca sử dụng uyệt đơn đăng ký}
	
	\paragraph{Biểu đồ ca sử dụng}
	\noindent
	
	\begin{figure}[H]
		\centering
		\includegraphics[width=0.9\linewidth]{images/uc_uc06}
		\caption{Biểu đồ ca sử dụng duyệt đơn đăng ký}
		\label{fig:ucuc06}
	\end{figure}
	
	
	\begin{center}
		\begin{tabular}{|p{4cm}|p{10cm}|}
			\hline
			\textbf{Tên ca sử dụng} & \textbf{Duyệt đơn đăng ký nội trú} \\
			\hline
			Tác nhân chính & Nhân viên quản lý \\
			\hline
			Dòng sự kiện chính &
			\begin{minipage}[t]{10cm}
				\vspace{2pt}
				\begin{enumerate}
					\item Nhân viên xem danh sách đơn đăng ký có trạng thái PENDING.
					\item Chọn đơn và xem chi tiết thông tin sinh viên và phòng đăng ký.
					\item Nhân viên ra quyết định phê duyệt hoặc từ chối.
					\item Nếu phê duyệt: cập nhật trạng thái đơn thành APPROVED, tăng \texttt{current\_occupancy} của phòng, cập nhật trạng thái phòng nếu đầy và gửi thông báo cho sinh viên.
					\item Nếu từ chối: cập nhật trạng thái đơn thành REJECTED, gửi thông báo từ chối cho sinh viên.
				\end{enumerate}
				\vspace{2pt}
			\end{minipage} \\
			\hline
			Dòng sự kiện phụ & Nếu tại thời điểm phê duyệt phòng đã đầy hoặc đang bảo trì, hệ thống cảnh báo và yêu cầu xác nhận lại. \\
			\hline
			Điều kiện tiên quyết & Nhân viên đã đăng nhập; tồn tại đơn đăng ký ở trạng thái PENDING. \\
			\hline
			Điều kiện sau & Đơn được cập nhật trạng thái; nếu APPROVED thì số người ở trong phòng tăng lên một và trạng thái phòng được rà soát lại. \\
			\hline
		\end{tabular}
	\end{center}
	
	\paragraph{Biểu đồ hoạt động}
	\noindent
	
	\begin{figure}[H]
		\centering
		\includegraphics[width=0.5\linewidth]{images/activity_uc06}
		\caption{Biểu đồ hoạt động duyệt đơn đăng ký}
		\label{fig:activityuc06}
	\end{figure}
	
	
	Nhân viên xem danh sách đơn PENDING và chọn đơn cần xử lý; nếu quyết định từ chối, cập nhật trạng thái thành REJECTED và gửi thông báo. Nếu đồng ý, hệ thống kiểm tra lại phòng còn hợp lệ trước khi cập nhật thành APPROVED, tăng số người ở thực tế, cập nhật trạng thái phòng nếu đạt sức chứa tối đa và gửi thông báo xác nhận.
	
	\paragraph{Biểu đồ tuần tự}
	\noindent
	
	Nhân viên gửi \texttt{PUT /api/v1/registrations/\{id\}/approve} hoặc \texttt{/reject}; Backend truy vấn trạng thái phòng trong Database, nếu phê duyệt thì \texttt{UPDATE} trạng thái đơn thành APPROVED, \texttt{UPDATE} \texttt{current\_occupancy} phòng và \texttt{INSERT} thông báo cho sinh viên trong cùng một giao dịch.
	
	\begin{figure}[H]
		\centering
		\includegraphics[width=1\linewidth]{images/sequence_uc06}
		\caption{Biểu đồ tuần tự duyệt đơn đăng ký}
		\label{fig:sequenceuc06}
	\end{figure}
	
	
	% -----------------------------------------------------------
	\subsubsection*{1.7 Ca sử dụng trả phòng}
	\addcontentsline{toc}{subsubsection}{1.7 Ca sử dụng trả phòng}
	
	\paragraph{Biểu đồ ca sử dụng}
	\noindent
	
	\begin{figure}[H]
		\centering
		\includegraphics[width=0.8\linewidth]{images/uc_uc07}
		\caption{Biểu đồ ca sử dụng trả phòng}
		\label{fig:ucuc07}
	\end{figure}
	
	
	\begin{center}
		\begin{tabular}{|p{4cm}|p{10cm}|}
			\hline
			\textbf{Tên ca sử dụng} & \textbf{Xử lý trả phòng} \\
			\hline
			Tác nhân chính & Nhân viên quản lý \\
			\hline
			Dòng sự kiện chính &
			\begin{minipage}[t]{10cm}
				\vspace{2pt}
				\begin{enumerate}
					\item Nhân viên tìm kiếm đăng ký APPROVED của sinh viên cần trả phòng.
					\item Xem thông tin chi tiết đăng ký và tình trạng hóa đơn của sinh viên.
					\item Xác nhận thực hiện trả phòng.
					\item Hệ thống cập nhật trạng thái đơn thành CHECKED\_OUT, ghi nhận ngày trả phòng thực tế.
					\item Giảm \texttt{current\_occupancy} của phòng và cập nhật trạng thái phòng nếu cần.
					\item Gửi thông báo xác nhận trả phòng cho sinh viên.
				\end{enumerate}
				\vspace{2pt}
			\end{minipage} \\
			\hline
			Dòng sự kiện phụ & Nếu sinh viên còn hóa đơn chưa thanh toán, hệ thống cảnh báo nhưng vẫn cho phép nhân viên tiếp tục thực hiện trả phòng. \\
			\hline
			Điều kiện tiên quyết & Tồn tại đăng ký ở trạng thái APPROVED của sinh viên cần trả phòng. \\
			\hline
			Điều kiện sau & Đơn đăng ký chuyển sang CHECKED\_OUT; số người ở trong phòng giảm đi một và trạng thái phòng được rà soát lại. \\
			\hline
		\end{tabular}
	\end{center}
	
	\paragraph{Biểu đồ hoạt động}
	\noindent
	
	\begin{figure}[H]
		\centering
		\includegraphics[width=0.6\linewidth]{images/activity_uc07}
		\caption{Biểu đồ hoạt động trả phòng}
		\label{fig:activityuc07}
	\end{figure}
	
	
	Nhân viên tìm và chọn đăng ký APPROVED cần xử lý; hệ thống kiểm tra và cảnh báo nếu sinh viên còn nợ hóa đơn. Sau khi xác nhận, hệ thống cập nhật trạng thái thành CHECKED\_OUT, giảm số người ở thực tế của phòng, cập nhật trạng thái phòng sang AVAILABLE nếu phòng vừa trống chỗ và gửi thông báo.
	
	\paragraph{Biểu đồ tuần tự}
	\noindent
	
	Nhân viên gửi \texttt{PUT /api/v1/registrations/\{id\}/checkout}; Backend \texttt{UPDATE} trạng thái đăng ký thành CHECKED\_OUT với \texttt{checkout\_date = now()}, \texttt{UPDATE} \texttt{current\_occupancy} phòng giảm một, rà soát và cập nhật trạng thái phòng, đồng thời \texttt{INSERT} thông báo cho sinh viên, tất cả trong cùng một giao dịch.
	
	\begin{figure}[H]
		\centering
		\includegraphics[width=1\linewidth]{images/sequence_uc07}
		\caption{Biểu đồ tuần tự trả phòng}
		\label{fig:sequenceuc07}
	\end{figure}
	
	
	
	% -----------------------------------------------------------
	\subsubsection*{1.8 Ca sử dụng tạo hóa đơn}
	\addcontentsline{toc}{subsubsection}{1.8 Ca sử dụng tạo hóa đơn}
	
	\paragraph{Biểu đồ ca sử dụng}
	\noindent
	
	\begin{figure}[H]
		\centering
		\includegraphics[width=0.85\linewidth]{images/uc_uc08}
		\caption{Biểu đồ ca sử dụng tạo hóa đơn}
		\label{fig:ucuc08}
	\end{figure}
	
	
	\begin{center}
		\begin{tabular}{|p{4cm}|p{10cm}|}
			\hline
			\textbf{Tên ca sử dụng} & \textbf{Tạo hóa đơn định kỳ} \\
			\hline
			Tác nhân chính & Nhân viên quản lý \\
			\hline
			Dòng sự kiện chính &
			\begin{minipage}[t]{10cm}
				\vspace{2pt}
				\begin{enumerate}
					\item Nhân viên chọn sinh viên đang cư trú hợp lệ và tháng/năm cần lập hóa đơn.
					\item Nhập chỉ số điện tiêu thụ (kWh) và nước tiêu thụ (m³).
					\item Hệ thống tự động tính: tổng = tiền phòng + điện $\times$ 3.500đ + nước $\times$ 15.000đ.
					\item Nhân viên kiểm tra và xác nhận tạo hóa đơn.
					\item Hệ thống lưu hóa đơn với trạng thái UNPAID và gửi thông báo cho sinh viên.
				\end{enumerate}
				\vspace{2pt}
			\end{minipage} \\
			\hline
			Dòng sự kiện phụ & Nếu hóa đơn cho sinh viên đó trong cùng tháng/năm đã tồn tại, hệ thống trả lỗi 409 Conflict và không tạo trùng. Nếu sinh viên không có đăng ký hợp lệ, hệ thống từ chối yêu cầu. \\
			\hline
			Điều kiện tiên quyết & Sinh viên có đăng ký với trạng thái APPROVED tại thời điểm lập hóa đơn; chưa tồn tại hóa đơn cùng tháng/năm cho cặp sinh viên -- phòng này. \\
			\hline
			Điều kiện sau & Hóa đơn mới được lưu với trạng thái UNPAID; sinh viên nhận thông báo và có thể xem chi tiết. \\
			\hline
		\end{tabular}
	\end{center}
	
	\paragraph{Biểu đồ hoạt động}
	\noindent
	
	Nhân viên chọn sinh viên và tháng/năm lập hóa đơn; hệ thống kiểm tra tránh trùng lặp hóa đơn trước khi tiếp nhận chỉ số điện và nước. Hệ thống tính tổng tiền theo công thức cố định, nhân viên xác nhận rồi hóa đơn được lưu với trạng thái UNPAID và sinh viên nhận thông báo.
	\begin{figure}[H]
		\centering
		\includegraphics[width=0.55\linewidth]{images/activity_uc08}
		\caption{Biểu đồ hoạt động tạo hóa đơn}
		\label{fig:activityuc08}
	\end{figure}
	
	
	
	\paragraph{Biểu đồ tuần tự}
	\noindent
	\begin{figure}[H]
		\centering
		\includegraphics[width=0.9\linewidth]{images/sequence_uc08}
		\caption{Biểu đồ tuần tự tạo hóa đơn}
		\label{fig:sequenceuc08}
	\end{figure}
	Frontend tính tạm thời tổng tiền và hiển thị cho nhân viên xem trước; sau khi xác nhận, gửi \texttt{POST /api/v1/invoices}; Backend kiểm tra trùng lặp (student + month + year + room) trong Database. Nếu hợp lệ thì \texttt{INSERT} hóa đơn với trạng thái UNPAID và \texttt{INSERT} thông báo, trả về 201 Created.
	
	
	
	% -----------------------------------------------------------
	\subsubsection*{1.9 Ca sử dụng thanh toán hóa đơn}
	\addcontentsline{toc}{subsubsection}{1.9 Ca sử dụng thanh toán hóa đơn}
	
	\paragraph{Biểu đồ ca sử dụng}
	\noindent
	
	\begin{figure}[H]
		\centering
		\includegraphics[width=0.75\linewidth]{images/uc_uc09}
		\caption{Biểu đồ ca sử dụng thanh toán hóa đơn}
		\label{fig:ucuc09}
	\end{figure}
	
	
	\begin{center}
		\begin{tabular}{|p{4cm}|p{10cm}|}
			\hline
			\textbf{Tên ca sử dụng} & \textbf{Xem và thanh toán hóa đơn} \\
			\hline
			Tác nhân chính & Sinh viên \\
			\hline
			Dòng sự kiện chính &
			\begin{minipage}[t]{10cm}
				\vspace{2pt}
				\begin{enumerate}
					\item Sinh viên vào mục Hóa đơn, hệ thống hiển thị danh sách hóa đơn theo tháng.
					\item Sinh viên chọn hóa đơn có trạng thái UNPAID để xem chi tiết.
					\item Xem tiền phòng, tiền điện, tiền nước và tổng cộng phải nộp.
					\item Sinh viên xác nhận thanh toán.
					\item Hệ thống cập nhật trạng thái hóa đơn thành PAID và ghi nhận thời điểm thanh toán.
				\end{enumerate}
				\vspace{2pt}
			\end{minipage} \\
			\hline
			Dòng sự kiện phụ & Nếu hóa đơn đã ở trạng thái PAID, hệ thống trả lỗi 400 Bad Request và không cho thanh toán lần thứ hai. \\
			\hline
			Điều kiện tiên quyết & Sinh viên đã đăng nhập; tồn tại hóa đơn ở trạng thái UNPAID thuộc về sinh viên đó. \\
			\hline
			Điều kiện sau & Hóa đơn chuyển sang trạng thái PAID với \texttt{paid\_at} ghi nhận thời điểm thanh toán. \\
			\hline
		\end{tabular}
	\end{center}
	
	\paragraph{Biểu đồ hoạt động}
	\noindent
	
	\begin{figure}[H]
		\centering
		\includegraphics[width=0.55\linewidth]{images/activity_uc09}
		\caption{Biểu đồ hoạt động thanh toán hóa đơn}
		\label{fig:activityuc09}
	\end{figure}
	
	
	Sinh viên vào trang Hóa đơn và chọn hóa đơn UNPAID để xem chi tiết; nếu xác nhận thanh toán, hệ thống cập nhật trạng thái thành PAID và ghi nhận ngày thanh toán rồi hiển thị xác nhận. Nếu hủy, sinh viên quay lại danh sách.
	
	\paragraph{Biểu đồ tuần tự}
	\noindent
	
	Sinh viên gửi \texttt{GET /api/v1/invoices/me} để lấy danh sách hóa đơn; sau đó gửi \texttt{PUT /api/v1/invoices/\{id\}/pay}; Backend kiểm tra trạng thái hóa đơn trong Database, nếu là UNPAID thì \texttt{UPDATE} sang PAID với \texttt{paid\_at = now()} và trả về 200 OK.
	\begin{figure}[H]
		\centering
		\includegraphics[width=0.9\linewidth]{images/sequence_uc09}
		\caption{Biểu đồ tuần tự thanh toán hóa đơn}
		\label{fig:sequenceuc09}
	\end{figure}
	
	
	
	% -----------------------------------------------------------
	\subsubsection*{1.10 Ca sử dụng yêu cầu bảo trì}
	\addcontentsline{toc}{subsubsection}{1.10 Ca sử dụng yêu cầu bảo trì}
	
	\paragraph{Biểu đồ ca sử dụng}
	\noindent
	
	\begin{figure}[H]
		\centering
		\includegraphics[width=0.9\linewidth]{images/uc_uc10}
		\caption{Biểu đồ ca sử dụng yêu cầu bảo trì}
		\label{fig:ucuc10}
	\end{figure}
	
	
	\begin{center}
		\begin{tabular}{|p{4cm}|p{10cm}|}
			\hline
			\textbf{Tên ca sử dụng} & \textbf{Gửi yêu cầu bảo trì} \\
			\hline
			Tác nhân chính & Sinh viên \\
			\hline
			Dòng sự kiện chính &
			\begin{minipage}[t]{10cm}
				\vspace{2pt}
				\begin{enumerate}
					\item Sinh viên vào mục Yêu cầu bảo trì.
					\item Hệ thống kiểm tra sinh viên có đăng ký APPROVED hợp lệ.
					\item Sinh viên nhập tiêu đề và mô tả chi tiết sự cố.
					\item Gửi yêu cầu; hệ thống tạo bản ghi với trạng thái PENDING và thông báo tiếp nhận.
					\item Sinh viên có thể theo dõi trạng thái xử lý trên cùng giao diện.
				\end{enumerate}
				\vspace{2pt}
			\end{minipage} \\
			\hline
			Dòng sự kiện phụ & Nếu sinh viên không có đăng ký APPROVED tại thời điểm gửi, hệ thống từ chối với lỗi 403 Forbidden. \\
			\hline
			Điều kiện tiên quyết & Sinh viên đã đăng nhập và có đăng ký phòng với trạng thái APPROVED. \\
			\hline
			Điều kiện sau & Yêu cầu bảo trì được tạo với trạng thái PENDING và nhân viên có thể thấy trong danh sách cần xử lý. \\
			\hline
		\end{tabular}
	\end{center}
	
	\paragraph{Biểu đồ hoạt động}
	\noindent
	
	Sinh viên vào mục Yêu cầu bảo trì; hệ thống kiểm tra điều kiện cư trú hợp lệ trước khi cho phép nhập thông tin. Sau khi điền tiêu đề và mô tả sự cố, sinh viên gửi yêu cầu và nhận thông báo tiếp nhận; trạng thái yêu cầu được đặt thành PENDING.	
	\begin{figure}[H]
		\centering
		\includegraphics[width=0.45\linewidth]{images/activity_uc10}
		\caption{Biểu đồ hoạt động yêu cầu bảo trì}
		\label{fig:activityuc10}
	\end{figure}
	
	
	
	\paragraph{Biểu đồ tuần tự}
	\noindent
	
	\begin{figure}[H]
		\centering
		\includegraphics[width=0.9\linewidth]{images/sequence_uc10}
		\caption{Biểu đồ tuần tự yêu cầu bảo trì}
		\label{fig:sequenceuc10}
	\end{figure}
	
	
	Sinh viên gửi \texttt{POST /api/v1/maintenance-requests}; Backend kiểm tra sinh viên có bản ghi trong \texttt{room\_registrations} với trạng thái APPROVED trong Database. Nếu hợp lệ thì \texttt{INSERT} yêu cầu bảo trì với trạng thái PENDING và trả về 201 Created; ngược lại trả 403 Forbidden.
	
	% -----------------------------------------------------------
	\subsubsection*{1.11 Ca sử dụng xử lý bảo trì}
	\addcontentsline{toc}{subsubsection}{1.11 Ca sử dụng xử lý bảo trì}
	
	\paragraph{Biểu đồ ca sử dụng}
	\noindent
	
	\begin{figure}[H]
		\centering
		\includegraphics[width=0.8\linewidth]{images/uc_uc11}
		\caption{Biểu đồ ca sử dụng xử lý bảo trì}
		\label{fig:ucuc11}
	\end{figure}
	
	
	\begin{center}
		\begin{tabular}{|p{4cm}|p{10cm}|}
			\hline
			\textbf{Tên ca sử dụng} & \textbf{Xử lý yêu cầu bảo trì} \\
			\hline
			Tác nhân chính & Nhân viên quản lý \\
			\hline
			Dòng sự kiện chính &
			\begin{minipage}[t]{10cm}
				\vspace{2pt}
				\begin{enumerate}
					\item Nhân viên xem danh sách yêu cầu bảo trì và lọc theo trạng thái.
					\item Chọn yêu cầu cần xử lý và xem chi tiết sự cố, phòng và sinh viên.
					\item Cập nhật trạng thái thành IN\_PROGRESS khi bắt đầu xử lý; gửi thông báo cho sinh viên.
					\item Sau khi hoàn thành sửa chữa, cập nhật trạng thái thành COMPLETED với ghi chú kết quả.
					\item Gửi thông báo hoàn thành cho sinh viên.
				\end{enumerate}
				\vspace{2pt}
			\end{minipage} \\
			\hline
			Dòng sự kiện phụ & Nhân viên có thể trực tiếp chuyển từ PENDING sang COMPLETED nếu vấn đề được giải quyết ngay. \\
			\hline
			Điều kiện tiên quyết & Nhân viên đã đăng nhập; tồn tại yêu cầu bảo trì ở trạng thái PENDING hoặc IN\_PROGRESS. \\
			\hline
			Điều kiện sau & Yêu cầu bảo trì được cập nhật trạng thái; sinh viên nhận thông báo về tiến độ xử lý. \\
			\hline
		\end{tabular}
	\end{center}
	
	\paragraph{Biểu đồ hoạt động}
	\noindent
	
	\begin{figure}[H]
		\centering
		\includegraphics[width=0.55\linewidth]{images/activity_uc11}
		\caption{Biểu đồ hoạt động xử lý bào trì}
		\label{fig:activityuc11}
	\end{figure}
	Nhân viên xem và lọc danh sách yêu cầu bảo trì; chọn yêu cầu và cập nhật trạng thái theo tiến trình xử lý: PENDING $\to$ IN\_PROGRESS $\to$ COMPLETED. Mỗi lần chuyển trạng thái đều kèm thông báo gửi đến sinh viên liên quan.
	
	\paragraph{Biểu đồ tuần tự}
	\noindent
	
	\begin{figure}[H]
		\centering
		\includegraphics[width=0.9\linewidth]{images/sequence_uc11}
		\caption{Biểu đồ tuần tự xử lý bảo trì}
		\label{fig:sequenceuc11}
	\end{figure}
	
	
	Nhân viên gửi \texttt{GET /api/v1/maintenance-requests} để lấy danh sách, sau đó gửi \texttt{PUT /api/v1/maintenance-requests/\{id\}} với trạng thái mới (IN\_PROGRESS hoặc COMPLETED) và ghi chú xử lý; Backend \texttt{UPDATE} bản ghi và \texttt{INSERT} thông báo cho sinh viên, trả về 200 OK.
	
	% -----------------------------------------------------------
	\subsubsection*{1.12 Ca sử dụng gửi thông báo}
	\addcontentsline{toc}{subsubsection}{1.12 ca sử dụng gửi thông báo}
	
	\paragraph{Biểu đồ ca sử dụng}
	\noindent
	
	\begin{figure}[H]
		\centering
		\includegraphics[width=0.7\linewidth]{images/uc_uc12}
		\caption{Biểu đồ ca sử dụng gửi thông báo}
		\label{fig:ucuc12}
	\end{figure}
	
	\begin{center}
		\begin{tabular}{|p{4cm}|p{10cm}|}
			\hline
			\textbf{Tên ca sử dụng} & \textbf{Soạn và gửi thông báo} \\
			\hline
			Tác nhân chính & Nhân viên quản lý, Quản trị viên \\
			\hline
			Dòng sự kiện chính &
			\begin{minipage}[t]{10cm}
				\vspace{2pt}
				\begin{enumerate}
					\item Nhân viên hoặc Quản trị viên soạn tiêu đề và nội dung thông báo.
					\item Chọn đối tượng nhận: ALL (tất cả), STUDENT (sinh viên) hoặc STAFF (nhân viên).
					\item Phát hành thông báo; hệ thống lưu với \texttt{target\_role} tương ứng.
					\item Người dùng thuộc đối tượng nhận thấy thông báo khi đăng nhập hoặc tải trang.
				\end{enumerate}
				\vspace{2pt}
			\end{minipage} \\
			\hline
			Dòng sự kiện phụ & Nếu nội dung thông báo rỗng hoặc thiếu trường bắt buộc, hệ thống hiển thị lỗi xác thực và không phát hành. \\
			\hline
			Điều kiện tiên quyết & Người thực hiện đã đăng nhập với vai trò STAFF hoặc ADMIN. \\
			\hline
			Điều kiện sau & Thông báo được lưu trong hệ thống; người dùng thuộc \texttt{target\_role} có thể xem khi truy cập trang thông báo. \\
			\hline
		\end{tabular}
	\end{center}
	
	\paragraph{Biểu đồ hoạt động}
	\noindent
	
	\begin{figure}[H]
		\centering
		\includegraphics[width=0.55\linewidth]{images/activity_uc12}
		\caption{Biểu đồ hoạt động gửi thông báo}
		\label{fig:activityuc12}
	\end{figure}
	
	Người gửi soạn tiêu đề, nội dung và chọn đối tượng nhận; hệ thống kiểm tra dữ liệu hợp lệ trước khi phát hành. Sau khi lưu thông báo vào cơ sở dữ liệu với \texttt{target\_role} phù hợp, người dùng thuộc đúng đối tượng sẽ thấy thông báo khi truy cập hệ thống.
	
	\paragraph{Biểu đồ tuần tự}
	\noindent
	
	\begin{figure}[H]
		\centering
		\includegraphics[width=0.9\linewidth]{images/sequence_uc12}
		\caption{Biểu đồ tuần tự gửi thông báo}
		\label{fig:sequenceuc12}
	\end{figure}
	
	Người gửi gửi \texttt{POST /api/v1/notifications} với tiêu đề, nội dung và \texttt{target\_role}; Backend kiểm tra quyền (ADMIN hoặc STAFF) và \texttt{INSERT} thông báo vào Database. Người nhận gửi \texttt{GET /api/v1/notifications}; Backend truy vấn và lọc theo \texttt{target\_role} phù hợp với vai trò người dùng hiện tại rồi trả về danh sách.
	
	% -----------------------------------------------------------
	\subsubsection*{1.13 Ca sử dụng xem báo cáo thống kê}
	\addcontentsline{toc}{subsubsection}{1.13 Ca sử dụng xem báo cáo thống kê}
	
	\paragraph{Biểu đồ ca sử dụng}
	\noindent
	
	\begin{figure}[H]
		\centering
		\includegraphics[width=0.9\linewidth]{images/uc_uc13}
		\caption{Biểu đồ ca sử dụng xem báo cáo thống kê}
		\label{fig:ucuc13}
	\end{figure}
	
	
	\begin{center}
		\begin{tabular}{|p{4cm}|p{10cm}|}
			\hline
			\textbf{Tên ca sử dụng} & \textbf{Xem báo cáo thống kê vận hành} \\
			\hline
			Tác nhân chính & Quản trị viên \\
			\hline
			Dòng sự kiện chính &
			\begin{minipage}[t]{10cm}
				\vspace{2pt}
				\begin{enumerate}
					\item Quản trị viên truy cập trang Báo cáo thống kê.
					\item Chọn loại báo cáo: doanh thu theo tháng hoặc tỷ lệ lấp đầy theo tòa nhà.
					\item Thiết lập tham số lọc (tháng/năm, tòa nhà).
					\item Hệ thống tổng hợp dữ liệu từ cơ sở dữ liệu và trả về kết quả.
					\item Hiển thị số liệu dưới dạng bảng hoặc biểu đồ trực quan.
				\end{enumerate}
				\vspace{2pt}
			\end{minipage} \\
			\hline
			Dòng sự kiện phụ & Nếu không có dữ liệu trong khoảng thời gian được chọn, hệ thống hiển thị kết quả rỗng với thông báo phù hợp. \\
			\hline
			Điều kiện tiên quyết & Người thực hiện đã đăng nhập với vai trò Quản trị viên. \\
			\hline
			Điều kiện sau & Số liệu thống kê được hiển thị; không có thay đổi dữ liệu trong hệ thống. \\
			\hline
		\end{tabular}
	\end{center}
	
	\paragraph{Biểu đồ hoạt động}
	\noindent
	
	\begin{figure}[H]
		\centering
		\includegraphics[width=0.65\linewidth]{images/activity_uc13}
		\caption{Biểu đồ hoạt động xem báo cáo thống kê}
		\label{fig:activityuc13}
	\end{figure}
	
	Quản trị viên truy cập trang Báo cáo và chọn loại báo cáo; hệ thống xử lý song song hai luồng: báo cáo doanh thu (tổng hợp \texttt{SUM(total\_amount)} từ hóa đơn PAID theo tháng/năm) và báo cáo tỷ lệ lấp đầy (tính \texttt{current\_occupancy / capacity} theo từng tòa nhà). Kết quả được hiển thị dưới dạng bảng thống kê.
	
	\paragraph{Biểu đồ tuần tự}
	\noindent
	
	\begin{figure}[H]
		\centering
		\includegraphics[width=0.9\linewidth]{images/sequence_uc13}
		\caption{Biểu đồ tuần tự xem báo cáo thống kê}
		\label{fig:sequenceuc13}
	\end{figure}
	
	Quản trị viên gửi \texttt{GET /api/v1/reports/revenue?month=X\&year=Y} hoặc \texttt{GET /api/v1/reports/occupancy?building\_id=...}; Backend thực hiện truy vấn tổng hợp trên Database (\texttt{SUM} hoặc tính tỷ lệ), xử lý kết quả và trả về dữ liệu cho Frontend hiển thị.
	
	% -----------------------------------------------------------
	\subsection*{2. Cơ sở dữ liệu}
	\addcontentsline{toc}{subsection}{2. Cơ sở dữ liệu}
	
	\paragraph{Sơ đồ thiết kế cơ sở dữ liệu}
	\noindent
	
	\begin{figure}[H]
		\centering
		\includegraphics[width=1\linewidth]{images/db}
		\caption{Mô hình cơ sở dữ liệu}
		\label{fig:db}
	\end{figure}
	
	\paragraph{Danh sách các đối tượng}
	Hệ thống sử dụng 7 bảng chính để lưu trữ thông tin người dùng, tòa nhà, phòng ở, đăng ký nội trú, hóa đơn, yêu cầu bảo trì và thông báo nội bộ. Danh sách các bảng được trình bày trong Bảng~\ref{tab:doi-tuong}.
	
	\begin{table}[H]
		\centering
		\caption{Các đối tượng dữ liệu chính}
		\label{tab:doi-tuong}
		\begin{tabular}{|p{1.5cm}|p{4cm}|p{8.5cm}|}
			\hline
			\centering\textbf{STT} & \centering\textbf{Tên bảng} & \centering\arraybackslash\textbf{Mô tả} \\
			\hline
			1 & users & Lưu thông tin tài khoản của tất cả người dùng hệ thống (sinh viên, nhân viên, quản trị viên). \\
			\hline
			2 & buildings & Lưu thông tin các tòa nhà trong khu ký túc xá. \\
			\hline
			3 & rooms & Lưu thông tin phòng ở theo từng tòa nhà, bao gồm sức chứa, loại phòng, giá và trạng thái. \\
			\hline
			4 & room\_registrations & Lưu thông tin đơn đăng ký nội trú của sinh viên, theo dõi trạng thái từ PENDING đến CHECKED\_OUT. \\
			\hline
			5 & invoices & Lưu hóa đơn định kỳ theo tháng, bao gồm chi phí phòng, điện, nước và trạng thái thanh toán. \\
			\hline
			6 & maintenance\_requests & Lưu yêu cầu bảo trì của sinh viên và trạng thái xử lý của nhân viên. \\
			\hline
			7 & notifications & Lưu thông báo nội bộ phân loại theo đối tượng nhận (ALL, STUDENT, STAFF). \\
			\hline
		\end{tabular}
	\end{table}
	
	\begin{table}[H]
		\centering
		\caption{Bảng users}
		\label{tab:users}
		\begin{tabular}{|p{1.5cm}|p{4cm}|p{4cm}|p{5.5cm}|}
			\hline
			\centering\textbf{STT} & \centering\textbf{Tên thuộc tính} & \centering\textbf{Kiểu dữ liệu} & \centering\arraybackslash\textbf{Mô tả} \\
			\hline
			1 & id & UUID & Mã định danh duy nhất của người dùng (khóa chính). \\
			\hline
			2 & full\_name & VARCHAR(255) & Họ và tên đầy đủ của người dùng. \\
			\hline
			3 & student\_code & VARCHAR(32) & Mã số sinh viên hoặc nhân viên, duy nhất trong hệ thống. \\
			\hline
			4 & email & VARCHAR(255) & Địa chỉ email đăng nhập, duy nhất. \\
			\hline
			5 & password\_hash & VARCHAR(255) & Mật khẩu đã được băm bằng bcrypt. \\
			\hline
			6 & phone & VARCHAR(20) & Số điện thoại liên hệ (tùy chọn). \\
			\hline
			7 & role & ENUM & Vai trò: STUDENT, STAFF hoặc ADMIN. \\
			\hline
			8 & gender & ENUM & Giới tính: MALE hoặc FEMALE. \\
			\hline
			9 & nationality & ENUM & Quốc tịch: VIETNAM hoặc LAOS. \\
			\hline
			10 & created\_at & TIMESTAMPTZ & Thời điểm tạo tài khoản. \\
			\hline
		\end{tabular}
	\end{table}
	
	\begin{table}[H]
		\centering
		\caption{Bảng buildings}
		\label{tab:buildings}
		\begin{tabular}{|p{1.5cm}|p{4cm}|p{4cm}|p{5.5cm}|}
			\hline
			\centering\textbf{STT} & \centering\textbf{Tên thuộc tính} & \centering\textbf{Kiểu dữ liệu} & \centering\arraybackslash\textbf{Mô tả} \\
			\hline
			1 & id & INT & Mã tòa nhà tự tăng (khóa chính). \\
			\hline
			2 & name & ENUM & Mã ký hiệu tòa nhà (K1, K2, K3, \ldots). \\
			\hline
			3 & total\_floors & INT & Tổng số tầng của tòa nhà. \\
			\hline
			4 & description & TEXT & Mô tả thêm về tòa nhà (tùy chọn). \\
			\hline
			5 & status & ENUM & Trạng thái hoạt động: ACTIVE hoặc INACTIVE. \\
			\hline
		\end{tabular}
	\end{table}
	
	\begin{table}[H]
		\centering
		\caption{Bảng rooms}
		\label{tab:rooms}
		\begin{tabular}{|p{1.5cm}|p{4cm}|p{4cm}|p{5.5cm}|}
			\hline
			\centering\textbf{STT} & \centering\textbf{Tên thuộc tính} & \centering\textbf{Kiểu dữ liệu} & \centering\arraybackslash\textbf{Mô tả} \\
			\hline
			1 & id & INT & Mã phòng tự tăng (khóa chính). \\
			\hline
			2 & building\_id & INT & Khóa ngoại tham chiếu bảng buildings. \\
			\hline
			3 & room\_number & VARCHAR(20) & Số phòng, duy nhất trong mỗi tòa nhà. \\
			\hline
			4 & floor & INT & Tầng của phòng. \\
			\hline
			5 & capacity & INT & Sức chứa tối đa của phòng. \\
			\hline
			6 & current\_occupancy & INT & Số sinh viên đang ở thực tế (mặc định 0). \\
			\hline
			7 & room\_type & ENUM & Loại phòng: MALE, FEMALE hoặc LAOS\_STUDENT. \\
			\hline
			8 & price\_per\_month & NUMERIC(12,2) & Giá phòng mỗi tháng (đồng). \\
			\hline
			9 & status & ENUM & Trạng thái: AVAILABLE, FULL hoặc MAINTENANCE. \\
			\hline
		\end{tabular}
	\end{table}
	
	\begin{table}[H]
		\centering
		\caption{Bảng room\_registrations}
		\label{tab:roomregistrations}
		\begin{tabular}{|p{1.5cm}|p{4cm}|p{4cm}|p{5.5cm}|}
			\hline
			\centering\textbf{STT} & \centering\textbf{Tên thuộc tính} & \centering\textbf{Kiểu dữ liệu} & \centering\arraybackslash\textbf{Mô tả} \\
			\hline
			1 & id & INT & Mã đăng ký tự tăng (khóa chính). \\
			\hline
			2 & student\_id & UUID & Khóa ngoại tham chiếu bảng users. \\
			\hline
			3 & room\_id & INT & Khóa ngoại tham chiếu bảng rooms. \\
			\hline
			4 & start\_date & DATE & Ngày bắt đầu ở dự kiến. \\
			\hline
			5 & end\_date & DATE & Ngày kết thúc ở dự kiến. \\
			\hline
			6 & status & ENUM & Trạng thái: PENDING, APPROVED, REJECTED hoặc CHECKED\_OUT. \\
			\hline
			7 & created\_at & TIMESTAMPTZ & Thời điểm tạo đơn đăng ký. \\
			\hline
		\end{tabular}
	\end{table}
	
	\begin{table}[H]
		\centering
		\caption{Bảng invoices}
		\label{tab:invoices}
		\begin{tabular}{|p{1.5cm}|p{4cm}|p{4cm}|p{5.5cm}|}
			\hline
			\centering\textbf{STT} & \centering\textbf{Tên thuộc tính} & \centering\textbf{Kiểu dữ liệu} & \centering\arraybackslash\textbf{Mô tả} \\
			\hline
			1 & id & INT & Mã hóa đơn tự tăng (khóa chính). \\
			\hline
			2 & student\_id & UUID & Khóa ngoại tham chiếu bảng users. \\
			\hline
			3 & room\_id & INT & Khóa ngoại tham chiếu bảng rooms. \\
			\hline
			4 & month & INT & Tháng lập hóa đơn (1--12). \\
			\hline
			5 & year & INT & Năm lập hóa đơn. \\
			\hline
			6 & electricity\_used\_kwh & NUMERIC(12,2) & Điện tiêu thụ tháng này (kWh). \\
			\hline
			7 & water\_used\_m3 & NUMERIC(12,2) & Nước tiêu thụ tháng này (m³). \\
			\hline
			8 & room\_fee & NUMERIC(12,2) & Tiền phòng theo tháng. \\
			\hline
			9 & electricity\_fee & NUMERIC(12,2) & Tiền điện = điện\_kWh $\times$ 3.500đ. \\
			\hline
			10 & water\_fee & NUMERIC(12,2) & Tiền nước = nước\_m³ $\times$ 15.000đ. \\
			\hline
			11 & total\_amount & NUMERIC(12,2) & Tổng tiền phải nộp. \\
			\hline
			12 & status & ENUM & Trạng thái: UNPAID hoặc PAID. \\
			\hline
			13 & due\_date & DATE & Hạn cuối thanh toán. \\
			\hline
			14 & paid\_at & TIMESTAMPTZ & Thời điểm thanh toán thực tế (NULL nếu chưa thanh toán). \\
			\hline
		\end{tabular}
	\end{table}
	
	\begin{table}[H]
		\centering
		\caption{Bảng maintenance\_requests}
		\label{tab:maintenancerequests}
		\begin{tabular}{|p{1.5cm}|p{4cm}|p{4cm}|p{5.5cm}|}
			\hline
			\centering\textbf{STT} & \centering\textbf{Tên thuộc tính} & \centering\textbf{Kiểu dữ liệu} & \centering\arraybackslash\textbf{Mô tả} \\
			\hline
			1 & id & INT & Mã yêu cầu tự tăng (khóa chính). \\
			\hline
			2 & student\_id & UUID & Khóa ngoại tham chiếu bảng users. \\
			\hline
			3 & room\_id & INT & Khóa ngoại tham chiếu bảng rooms. \\
			\hline
			4 & title & VARCHAR(255) & Tiêu đề mô tả ngắn về sự cố. \\
			\hline
			5 & description & TEXT & Mô tả chi tiết nội dung yêu cầu. \\
			\hline
			6 & status & ENUM & Trạng thái: PENDING, IN\_PROGRESS hoặc COMPLETED. \\
			\hline
			7 & created\_at & TIMESTAMPTZ & Thời điểm gửi yêu cầu. \\
			\hline
			8 & resolved\_at & TIMESTAMPTZ & Thời điểm hoàn thành xử lý (NULL nếu chưa xong). \\
			\hline
		\end{tabular}
	\end{table}
	
	\begin{table}[H]
		\centering
		\caption{Bảng notifications}
		\label{tab:notifications}
		\begin{tabular}{|p{1.5cm}|p{4cm}|p{4cm}|p{5.5cm}|}
			\hline
			\centering\textbf{STT} & \centering\textbf{Tên thuộc tính} & \centering\textbf{Kiểu dữ liệu} & \centering\arraybackslash\textbf{Mô tả} \\
			\hline
			1 & id & INT & Mã thông báo tự tăng (khóa chính). \\
			\hline
			2 & title & VARCHAR(255) & Tiêu đề thông báo. \\
			\hline
			3 & content & TEXT & Nội dung chi tiết thông báo. \\
			\hline
			4 & target\_role & ENUM & Đối tượng nhận: ALL, STUDENT hoặc STAFF. \\
			\hline
			5 & created\_by & UUID & Khóa ngoại tham chiếu bảng users (người tạo). \\
			\hline
			6 & created\_at & TIMESTAMPTZ & Thời điểm tạo thông báo. \\
			\hline
		\end{tabular}
	\end{table}
	
	\newpage
	\section*{CHƯƠNG 3: XÂY DỰNG CHƯƠNG TRÌNH}
	\addcontentsline{toc}{section}{CHƯƠNG 3: XÂY DỰNG CHƯƠNG TRÌNH}
	
	% -----------------------------------------------------------
	\subsection*{3.1. Kiến trúc hệ thống}
	\addcontentsline{toc}{subsection}{3.1. Kiến trúc hệ thống}
	
	\begin{figure}[H]
		\centering
		\begin{tikzpicture}[
			outerbox/.style={draw, thick, rounded corners=3pt,
				minimum width=8.5cm, minimum height=2cm},
			comp/.style={draw, dashed, rounded corners=2pt,
				minimum width=6.5cm, minimum height=0.9cm,
				text centered, font=\small},
			lbl/.style={font=\small\itshape}
			]
			%% Client
			\node[outerbox, fill=gray!12] at (0, 0) {};
			\node[font=\bfseries\small, anchor=north west]
			at (-4.15, 0.85) {Client};
			\node[comp] at (0, -0.1) {Trình duyệt web};
			
			%% Web Server
			\node[outerbox, fill=blue!8] at (0, -3.5) {};
			\node[font=\bfseries\small, anchor=north west]
			at (-4.15, -2.65) {Web Server};
			\node[comp] at (0, -3.6)
			{Ứng dụng React (Giao diện người dùng)};
			
			%% Application Server
			\node[outerbox, fill=green!8] at (0, -7) {};
			\node[font=\bfseries\small, anchor=north west]
			at (-4.15, -6.15) {Application Server};
			\node[comp] at (0, -7.1)
			{FastAPI --- Xử lý nghiệp vụ};
			
			%% Database Server
			\node[outerbox, fill=yellow!10] at (0, -10.5) {};
			\node[font=\bfseries\small, anchor=north west]
			at (-4.15, -9.65) {Database Server};
			\node[comp] at (0, -10.6)
			{Cơ sở dữ liệu PostgreSQL};
			
			%% Arrows (bottom of upper box → top of lower box, gap = 1.5cm)
			\draw[->, thick] (0, -1)   -- node[lbl, right=0.15cm] {HTTP/HTTPS}       (0, -2.5);
			\draw[->, thick] (0, -4.5) -- node[lbl, right=0.15cm] {REST API (JSON)}  (0, -6);
			\draw[->, thick] (0, -8)   -- node[lbl, right=0.15cm] {SQL / asyncpg}    (0, -9.5);
		\end{tikzpicture}
		\caption{Kiến trúc tổng thể của hệ thống quản lý ký túc xá}
		\label{fig:system-architecture}
	\end{figure}
	
	Hệ thống được xây dựng theo mô hình kiến trúc nhiều tầng (multi-tier architecture) với bốn lớp được phân tách rõ ràng theo chức năng. Tầng Client là trình duyệt web của người dùng -- nơi khởi phát mọi yêu cầu tương tác và nhận lại kết quả hiển thị từ phía máy chủ. Tầng Web Server đảm nhiệm việc phục vụ ứng dụng giao diện đơn trang (SPA) được xây dựng bằng React, cho phép người dùng tương tác liền mạch mà không cần tải lại toàn bộ trang. Tầng Application Server chứa toàn bộ logic nghiệp vụ, tiếp nhận yêu cầu từ giao diện qua giao thức REST API, thực hiện xác thực danh tính bằng cơ chế JWT, kiểm tra phân quyền và áp dụng các ràng buộc nghiệp vụ trước khi truy xuất hoặc thay đổi dữ liệu. Tầng Database Server chịu trách nhiệm lưu trữ và quản lý toàn bộ dữ liệu vận hành của hệ thống thông qua cơ sở dữ liệu quan hệ PostgreSQL, với mọi thao tác I/O được thực hiện theo phương thức bất đồng bộ nhằm tối ưu hiệu năng xử lý đồng thời. Toàn bộ bốn thành phần này được đóng gói và vận hành trong môi trường container hóa, đảm bảo tính nhất quán và khả năng tái triển khai giữa các môi trường phát triển và sản xuất thực tế.
	
	% -----------------------------------------------------------
	\subsection*{3.2. Công nghệ sử dụng}
	\addcontentsline{toc}{subsection}{3.2. Công nghệ sử dụng}
	
	\subsubsection*{3.2.1. Python và FastAPI}
	\addcontentsline{toc}{subsubsection}{3.2.1. Python và FastAPI}
	
	Phía máy chủ ứng dụng, hệ thống sử dụng Python kết hợp với framework FastAPI để xây dựng toàn bộ lớp API RESTful. FastAPI được lựa chọn nhờ khả năng xử lý yêu cầu bất đồng bộ thông qua cú pháp \texttt{async/await} của Python, giúp tối ưu hiệu năng khi có nhiều kết nối đồng thời từ phía client. Framework này tích hợp chặt chẽ với thư viện Pydantic để tự động kiểm tra và chuyển đổi kiểu dữ liệu đầu vào, đồng thời hỗ trợ cơ chế Dependency Injection giúp tổ chức code theo hướng mô-đun, tách biệt rõ ràng giữa các tầng xác thực, phân quyền và xử lý nghiệp vụ. Cơ chế xác thực danh tính được thực hiện qua JSON Web Token (JWT) với cặp Access Token và Refresh Token được lưu trong HttpOnly cookie, kết hợp bcrypt để băm mật khẩu người dùng nhằm đảm bảo an toàn thông tin xác thực ngay cả trong trường hợp cơ sở dữ liệu bị truy cập trái phép.
	
	\begin{figure}[H]
		\centering
		\includegraphics[width=0.75\linewidth]{images/fastapi}
		\caption{FastAPI -- Framework xây dựng API phía backend}
		\label{fig:fastapi}
	\end{figure}
	
	\subsubsection*{3.2.2. PostgreSQL, SQLAlchemy và Alembic}
	\addcontentsline{toc}{subsubsection}{3.2.2. PostgreSQL, SQLAlchemy và Alembic}
	
	PostgreSQL được chọn làm hệ quản trị cơ sở dữ liệu quan hệ cho toàn bộ hệ thống nhờ độ ổn định cao, hỗ trợ đầy đủ các kiểu dữ liệu cần thiết như UUID, ENUM, TIMESTAMPTZ và khả năng xử lý giao dịch (transaction) mạnh mẽ -- điều đặc biệt quan trọng khi nhiều thao tác nghiệp vụ yêu cầu cập nhật đồng thời nhiều bảng trong cùng một giao dịch nguyên tử. Toàn bộ tương tác với cơ sở dữ liệu được thực hiện thông qua SQLAlchemy 2.0 theo chế độ bất đồng bộ hoàn toàn (async), kết hợp driver \texttt{asyncpg} để đạt hiệu năng I/O tối ưu; các đối tượng nghiệp vụ như người dùng, phòng ở, đơn đăng ký hay hóa đơn được định nghĩa dưới dạng model ORM, giúp mã nguồn rõ ràng và dễ bảo trì hơn so với viết truy vấn SQL thuần. Alembic đảm nhận vai trò quản lý lịch sử thay đổi lược đồ cơ sở dữ liệu thông qua các file migration được đánh số phiên bản, cho phép theo dõi và áp dụng thay đổi cấu trúc một cách có kiểm soát trong suốt quá trình phát triển.
	
	\begin{figure}[H]
		\centering
		\includegraphics[width=0.75\linewidth]{images/postgres}
		\caption{PostgreSQL -- Hệ quản trị cơ sở dữ liệu quan hệ}
		\label{fig:tech-postgresql}
	\end{figure}
	
	\subsubsection*{3.2.3. React, TypeScript và Vite}
	\addcontentsline{toc}{subsubsection}{3.2.3. React, TypeScript và Vite}
	
	Phía giao diện, hệ thống được xây dựng dưới dạng ứng dụng đơn trang (Single Page Application) sử dụng React 19 kết hợp TypeScript. React cho phép tổ chức giao diện theo hướng component hóa -- mỗi thành phần giao diện được xây dựng độc lập, có thể tái sử dụng linh hoạt ở nhiều màn hình khác nhau, từ bảng danh sách phòng đến form nhập liệu hóa đơn. TypeScript bổ sung hệ thống kiểu tĩnh vào mã nguồn JavaScript, giúp phát hiện sớm các lỗi không tương thích kiểu dữ liệu ngay tại thời điểm viết code thay vì chờ đến khi chạy chương trình, từ đó nâng cao độ tin cậy của ứng dụng. Vite được sử dụng làm công cụ build và môi trường phát triển, cung cấp tốc độ khởi động và hot-reload nhanh hơn đáng kể so với các giải pháp truyền thống, giúp rút ngắn chu kỳ phát triển -- chỉnh sửa -- kiểm tra.
	
	\begin{figure}[H]
		\centering
		\includegraphics[width=0.75\linewidth]{images/react-vite}
		\caption{React và TypeScript -- Nền tảng xây dựng giao diện người dùng}
		\label{fig:react-vite}
	\end{figure}
	
	\subsubsection*{3.2.4. TanStack Router, TanStack Query và Zustand}
	\addcontentsline{toc}{subsubsection}{3.2.4. TanStack Router, TanStack Query và Zustand}
	
	Để quản lý điều hướng phía client, hệ thống sử dụng TanStack Router -- một thư viện định tuyến tích hợp sẵn cơ chế bảo vệ route (route guard), cho phép kiểm tra vai trò người dùng trước khi cho phép truy cập vào các trang dành riêng cho quản trị viên, nhân viên hay sinh viên. Việc đồng bộ dữ liệu giữa frontend và backend được xử lý bởi TanStack Query (React Query), thư viện này tự động quản lý vòng đời của các yêu cầu HTTP, bao gồm cơ chế cache, refetch khi dữ liệu cũ và cập nhật lạc quan (optimistic update) sau các thao tác ghi -- giúp giao diện phản hồi nhanh và nhất quán với trạng thái thực của dữ liệu trên máy chủ. Zustand được sử dụng để lưu trữ các trạng thái toàn cục phía client như thông tin phiên đăng nhập hiện tại và trạng thái các thông báo trong ứng dụng, với cú pháp gọn nhẹ không đòi hỏi boilerplate phức tạp như các giải pháp quản lý trạng thái truyền thống.
	
	\subsubsection*{3.2.5. Shadcn/UI, Tailwind CSS và Recharts}
	\addcontentsline{toc}{subsubsection}{3.2.5. Shadcn/UI, Tailwind CSS và Recharts}
	
	Hệ thống giao diện được xây dựng dựa trên Shadcn/UI -- một bộ component sẵn sàng sử dụng được tổ chức theo nguyên tắc Radix UI, cung cấp đầy đủ các thành phần như dialog, dropdown, bảng dữ liệu, form và toast notification với khả năng tiếp cận (accessibility) tích hợp sẵn. Tailwind CSS đóng vai trò nền tảng định dạng theo hướng utility-first, cho phép xây dựng và điều chỉnh giao diện trực tiếp trong JSX mà không cần viết file CSS riêng, đồng thời đảm bảo tính nhất quán về khoảng cách, màu sắc và kiểu chữ trên toàn bộ ứng dụng. Recharts được tích hợp để hiển thị dữ liệu thống kê trên trang báo cáo dưới dạng biểu đồ cột và biểu đồ đường, giúp quản trị viên trực quan hóa xu hướng doanh thu thu phí và tỷ lệ lấp đầy phòng theo từng tháng một cách dễ đọc.
	
	\begin{figure}[H]
		\centering
		\includegraphics[width=0.75\linewidth]{images/tailwind-css}
		\caption{Tailwind CSS và Shadcn/UI -- Hệ thống giao diện người dùng}
		\label{fig:tailwind-css}
	\end{figure}
	
	\subsubsection*{3.2.6. Docker Compose}
	\addcontentsline{toc}{subsubsection}{3.2.6. Docker Compose}
	
	Toàn bộ hệ thống được đóng gói và vận hành bằng Docker Compose với ba dịch vụ chính: backend (FastAPI), frontend (React/Nginx) và cơ sở dữ liệu (PostgreSQL). Mỗi dịch vụ chạy trong container riêng biệt với cấu hình môi trường được quản lý qua file \texttt{.env}, đảm bảo rằng mọi thành viên trong nhóm phát triển đều làm việc trên cùng một phiên bản phần mềm và cấu hình, loại bỏ các lỗi phát sinh do sự khác biệt môi trường. Cách tiếp cận này cũng tạo điều kiện thuận lợi cho việc triển khai lên môi trường sản xuất -- chỉ cần sao chép toàn bộ cấu hình Docker Compose sang máy chủ đích và khởi chạy, hệ thống sẽ hoạt động nhất quán mà không cần cấu hình thủ công từng thành phần.
	
	\begin{figure}[H]
		\centering
		\includegraphics[width=0.75\linewidth]{images/docker}
		\caption{Docker Compose -- Môi trường container hóa và triển khai hệ thống}
		\label{fig:docker}
	\end{figure}
	
	% -----------------------------------------------------------
	\subsection*{3.3. Giao diện của hệ thống}
	\addcontentsline{toc}{subsection}{3.3. Giao diện của hệ thống}
	
	\subsubsection*{3.3.1. Giao diện đăng nhập}
	\addcontentsline{toc}{subsubsection}{3.3.1. Giao diện đăng nhập}
	
	\begin{figure}[H]
		\centering
		\includegraphics[width=0.85\linewidth]{images/ui_login}
		\caption{Giao diện đăng nhập}
		\label{fig:uilogin}
	\end{figure}
	
	Giao diện đăng nhập là cửa ngõ đầu tiên để người dùng tiếp cận hệ thống, với hai trường nhập liệu email và mật khẩu cùng phản hồi lỗi tức thời khi thông tin xác thực không hợp lệ. Sau khi đăng nhập thành công, hệ thống tự động điều hướng người dùng đến trang tổng quan phù hợp với vai trò được gán (quản trị viên, nhân viên quản lý hoặc sinh viên).
	
	\subsubsection*{3.3.2. Trang tổng quan}
	\addcontentsline{toc}{subsubsection}{3.3.2. Trang tổng quan}
	
	\begin{figure}[H]
		\centering
		\includegraphics[width=0.85\linewidth]{images/ui_dashboard}
		\caption{Trang tổng quan dành cho quản trị viên}
		\label{fig:uidashboard}
	\end{figure}
	
	\begin{figure}[H]
		\centering
		\includegraphics[width=0.85\linewidth]{images/ui_dashboard_staff}
		\caption{Trang tổng quan dành cho nhân viên quản lý tòa nhà}
		\label{fig:uidashboardstaff}
	\end{figure}
	
	
	Trang tổng quan cung cấp cái nhìn nhanh về tình trạng vận hành của hệ thống, bao gồm số lượng phòng đang hoạt động, số đơn đăng ký đang chờ xử lý, số hóa đơn chưa thanh toán và số yêu cầu bảo trì đang mở. Các chỉ số này được tính toán từ cơ sở dữ liệu khi trang được tải, giúp quản trị viên và nhân viên quản lý tòa nhà nhanh chóng nắm bắt trạng thái vận hành tổng thể mà không cần tra cứu từng mục riêng lẻ.
	
	\subsubsection*{3.3.3. Quản lý tòa nhà và phòng ở}
	\addcontentsline{toc}{subsubsection}{3.3.3. Quản lý tòa nhà và phòng ở}
	
	\begin{figure}[H]
		\centering
		\includegraphics[width=0.85\linewidth]{images/ui_building_management}
		\caption{Giao diện quản lý tòa nhà}
		\label{fig:uibuildingmanagement}
	\end{figure}
	
		
	Giao diện này cho phép quản trị viên xem và quản lý danh sách tòa nhà cùng toàn bộ phòng ở trong từng khu, với đầy đủ thông tin về sức chứa, số người đang ở thực tế, loại phòng (nam/nữ/lưu học sinh Lào) và trạng thái hiện tại. Quản trị viên có thể thêm mới, chỉnh sửa thông tin hoặc đánh dấu phòng đang bảo trì trực tiếp từ giao diện danh sách.
	
	\begin{figure}[H]
		\centering
		\includegraphics[width=0.85\linewidth]{images/ui_room_management}
		\caption{Giao diện quản lý phòng ở}
		\label{fig:uiroommanagement}
	\end{figure}
	

	
	\subsubsection*{3.3.4. Quản lý đăng ký nội trú}
	\addcontentsline{toc}{subsubsection}{3.3.4. Quản lý đăng ký nội trú}
	
	\begin{figure}[H]
		\centering
		\includegraphics[width=0.85\linewidth]{images/ui_registration_management}
		\caption{Giao diện duyệt đơn đăng ký nội trú}
		\label{fig:uiregistrationmanagement}
	\end{figure}
	
	Giao diện duyệt đơn tập trung toàn bộ các đơn đăng ký đang chờ xử lý, hiển thị thông tin sinh viên, phòng được yêu cầu và thời điểm nộp đơn để nhân viên có đủ căn cứ ra quyết định. Khi phê duyệt, hệ thống tự động cập nhật số người ở thực tế trong phòng và gửi thông báo xác nhận đến sinh viên; khi từ chối, trạng thái đơn được cập nhật tức thì và sinh viên nhận thông báo lý do.
	
	\subsubsection*{3.3.5. Quản lý hóa đơn}
	\addcontentsline{toc}{subsubsection}{3.3.5. Quản lý hóa đơn}
	
	\begin{figure}[H]
		\centering
		\includegraphics[width=0.85\linewidth]{images/ui_invoice_management}
		\caption{Giao diện quản lý hóa đơn hàng tháng}
		\label{fig:uiinvoicemanagement}
	\end{figure}
	
	Giao diện quản lý hóa đơn cho phép nhân viên lập hóa đơn hàng tháng cho từng sinh viên bằng cách nhập chỉ số điện và nước; hệ thống tự động tính tổng tiền phải nộp và lưu hóa đơn với trạng thái chưa thanh toán. Nhân viên có thể theo dõi danh sách hóa đơn theo trạng thái và xác nhận thanh toán khi sinh viên nộp tiền trực tiếp.
	
	\subsubsection*{3.3.6. Quản lý yêu cầu bảo trì}
	\addcontentsline{toc}{subsubsection}{3.3.6. Quản lý yêu cầu bảo trì}
	
	\begin{figure}[H]
		\centering
		\includegraphics[width=0.85\linewidth]{images/ui_maitenance}
		\caption{Giao diện quản lý và xử lý yêu cầu bảo trì}
		\label{fig:uimaitenance}
	\end{figure}
	
	Giao diện hiển thị toàn bộ yêu cầu bảo trì từ sinh viên kèm thông tin chi tiết về phòng, mô tả sự cố và thời điểm gửi; nhân viên có thể lọc theo trạng thái (Đang chờ / Đang xử lý / Hoàn thành) để ưu tiên xử lý. Mỗi lần cập nhật trạng thái, hệ thống tự động gửi thông báo tiến độ đến sinh viên liên quan.
	
	\subsubsection*{3.3.7. Báo cáo thống kê}
	\addcontentsline{toc}{subsubsection}{3.3.7. Báo cáo thống kê}
	
	\begin{figure}[H]
		\centering
		\includegraphics[width=0.85\linewidth]{images/ui_reports}
		\caption{Giao diện báo cáo thống kê vận hành}
		\label{fig:uireports}
	\end{figure}
	
	
	Trang báo cáo cung cấp hai nhóm thống kê chính: doanh thu thu phí theo tháng (tổng hợp từ các hóa đơn đã thanh toán) và tỷ lệ lấp đầy theo từng tòa nhà. Dữ liệu được trình bày dưới dạng bảng và biểu đồ trực quan, cho phép quản trị viên theo dõi xu hướng vận hành và đưa ra các quyết định điều chỉnh kịp thời.
	
	\subsubsection*{3.3.8. Quản lý người dùng}
	\addcontentsline{toc}{subsubsection}{3.3.7. Quản lý người dùng}
	\begin{figure}[H]
		\centering
		\includegraphics[width=0.85\linewidth]{images/ui_user_management}
		\caption{Giao diện quản lý người dùng cho quản trị viện}
		\label{fig:uiusermanagement}
	\end{figure}
	
	\subsubsection*{3.3.9. Tạo và quản lý thông báo}
	\addcontentsline{toc}{subsubsection}{3.3.9. Tạo và quản lý thông báo}
	\begin{figure}[H]
		\centering
		\includegraphics[width=0.85\linewidth]{images/ui_reports2}
		\caption{Giao diện tạo thông báo}
		\label{fig:uireports2}
	\end{figure}
	
	
	\subsubsection*{3.3.10. Giao diện dành cho sinh viên}
	\addcontentsline{toc}{subsubsection}{3.3.10. Giao diện dành cho sinh viên}
	
	\begin{figure}[H]
		\centering
		\includegraphics[width=0.85\linewidth]{images/ui_dashboard_student}
		\caption{Giao diện sinh viên -- Giao diện trang tổng quan}
		\label{fig:uidashboardstudent}
	\end{figure}
	
	Giao diện dành cho sinh viên tập trung vào bốn nhóm chức năng: xem danh sách phòng còn chỗ và nộp đơn đăng ký; theo dõi trạng thái đơn đăng ký và lịch sử cư trú; kiểm tra và xác nhận thanh toán hóa đơn hàng tháng; gửi và theo dõi tiến độ xử lý yêu cầu bảo trì. Toàn bộ thao tác được thực hiện qua giao diện đơn giản, nhất quán và hoàn toàn bằng tiếng Việt, giảm thiểu rào cản sử dụng đối với sinh viên ít tiếp xúc với phần mềm quản lý.
	
	\newpage
	\section*{KẾT LUẬN}
	\addcontentsline{toc}{section}{KẾT LUẬN}
	
	Đề tài đã hoàn thành mục tiêu xây dựng một ứng dụng web hỗ trợ số hóa toàn bộ chu trình quản lý nội trú sinh viên, từ giai đoạn đăng ký phòng cho đến khi sinh viên chính thức trả phòng và thanh lý hợp đồng cư trú.
	
	Hệ thống đã triển khai đầy đủ các nhóm chức năng nghiệp vụ cốt lõi: xác thực và phân quyền theo ba vai trò (quản trị viên, nhân viên quản lý và sinh viên); quản lý danh mục tòa nhà và phòng ở bao gồm phân loại đối tượng cư trú; quy trình đăng ký -- xét duyệt -- thanh lý nội trú với cơ chế tự động cập nhật tình trạng phòng; lập và theo dõi hóa đơn hàng tháng cho từng sinh viên; tiếp nhận và xử lý yêu cầu bảo trì theo vòng đời trạng thái; gửi thông báo nội bộ đến đúng nhóm đối tượng; và tổng hợp báo cáo thống kê vận hành theo tháng. Xuyên suốt quá trình xây dựng, một yêu cầu thiết kế được duy trì nhất quán là mọi thao tác nghiệp vụ đều phải đi qua tầng xử lý trung tâm để kiểm tra điều kiện và thực thi ràng buộc, đảm bảo tính toàn vẹn dữ liệu ngay cả trong các tình huống đồng thời hay ngoại lệ nghiệp vụ.
	
	Tuy vậy, hệ thống trong phiên bản hiện tại vẫn còn một số điểm cần hoàn thiện. Cơ chế thanh toán hóa đơn hiện chỉ dừng ở mức xác nhận thủ công, chưa tích hợp cổng thanh toán trực tuyến để sinh viên có thể thực hiện giao dịch không tiếp xúc. Thông tin sinh viên chưa được đồng bộ tự động từ hệ thống quản lý đào tạo của nhà trường, dẫn đến việc nhập liệu thủ công tốn thời gian và tiềm ẩn nguy cơ sai sót. Cơ chế thông báo hiện hoạt động theo phương thức kéo -- người dùng chỉ nhận được nội dung mới khi chủ động mở trang -- trong khi một cơ chế đẩy thời gian thực sẽ nâng cao đáng kể trải nghiệm tương tác. Ngoài ra, quy trình lập hóa đơn yêu cầu nhân viên nhập tay từng bản ghi; một cơ chế tự động hóa theo chu kỳ tháng sẽ giúp giảm đáng kể khối lượng thao tác lặp lại trong vận hành thực tế.
	
	Những định hướng cải tiến trên đều nằm trong phạm vi khả thi về mặt nghiệp vụ và có giá trị ứng dụng thực tiễn cao, là cơ sở để hệ thống tiếp tục được phát triển và mở rộng nhằm đáp ứng trọn vẹn hơn nhu cầu vận hành thực tế tại khu ký túc xá Trường Đại học Tây Bắc.
	
	\newpage
	\section*{}
	\addcontentsline{toc}{section}{TÀI LIỆU THAM KHẢO}
	
	\begin{thebibliography}{99}
		\bibitem{utb_history}
		Trường Đại học Tây Bắc, \textit{Lịch sử hình thành}, trang chính thức Trường Đại học Tây Bắc. Truy cập tại: \url{https://utb.edu.vn/lich-su-hinh-thanh.html}
		
		\bibitem{utb_khaigiang2024}
		Báo Sơn La, \textit{Xứng đáng là trường đào tạo nguồn nhân lực chất lượng cho tỉnh Sơn La và khu vực Tây Bắc}, 2024. Truy cập tại: \url{https://baosonla.vn}
		
		\bibitem{utb_ktx}
		Trường Đại học Tây Bắc, \textit{Thông tin tuyển sinh -- Ký túc xá}, trang chính thức Trường Đại học Tây Bắc. Truy cập tại: \url{https://utb.edu.vn/tuyen-sinh.html}
		
		\bibitem{utb_phong_ktu}
		Navigates, \textit{Trường Đại học Tây Bắc (TBU) -- Mã trường TTB}, 2025. Truy cập tại: \url{https://navigates.vn/truong-hoc/dai-hoc-tay-bac/}
		
		\bibitem{utb_phongsban}
		Trường Đại học Tây Bắc, \textit{Các phòng ban -- Phòng Quản trị Thiết bị}, trang chính thức Trường Đại học Tây Bắc. Truy cập tại: \url{https://utb.edu.vn/cac-phong-ban.html}
		
		\bibitem{starrez}
		StarRez, \textit{Student Housing Software and Residential Community Management}, tài liệu giới thiệu sản phẩm và giải pháp quản lý cư trú sinh viên, truy cập từ trang chính thức của StarRez.
		
		\bibitem{erezlife}
		eRezLife, \textit{Housing and Residence Life Software}, tài liệu giới thiệu nền tảng quản lý ký túc xá và đời sống nội trú sinh viên, truy cập từ trang chính thức của eRezLife.
		
		\bibitem{fastapi}
		Sebastián Ramírez, \textit{FastAPI Documentation}, tài liệu chính thức của FastAPI về xây dựng API, dependency injection và xử lý xác thực.
		
		\bibitem{sqlalchemy}
		SQLAlchemy Authors, \textit{SQLAlchemy Documentation}, tài liệu chính thức về ORM, truy vấn dữ liệu và ánh xạ đối tượng trong Python.
		
		\bibitem{postgresql}
		PostgreSQL Global Development Group, \textit{PostgreSQL Documentation}, tài liệu chính thức về thiết kế và vận hành cơ sở dữ liệu PostgreSQL.
		
		\bibitem{react}
		Meta and React Contributors, \textit{React Documentation}, tài liệu chính thức về phát triển giao diện người dùng theo mô hình component.
		
		\bibitem{docker}
		Docker Inc., \textit{Docker Documentation}, tài liệu chính thức về container hóa và Docker Compose.
		
		\bibitem{jwt}
		M. Jones, J. Bradley, N. Sakimura, \textit{RFC 7519: JSON Web Token (JWT)}, Internet Engineering Task Force, 2015.
	\end{thebibliography}
\end{document}